%%=============================================================================
%% Algèbre Abstraite II — Corps et Théorie de Galois
%% Licence L3 — Mathématiques
%% Langue : français
%%=============================================================================
\documentclass[12pt,a4paper,openany]{book}

%%— Encodage et langue --------------------------------------------------------
\usepackage[utf8]{inputenc}
\usepackage[T1]{fontenc}
\usepackage[french]{babel}

%%— Mathématiques --------------------------------------------------------------
\usepackage{amsmath}
\usepackage{amssymb}
\usepackage{amsthm}
\usepackage{mathtools}
\usepackage{mathrsfs}

%%— Graphiques et diagrammes ---------------------------------------------------
\usepackage{tikz}
\usetikzlibrary{calc}
\usepackage{tikz-cd}

%%— Boîtes colorées -----------------------------------------------------------
\usepackage[most]{tcolorbox}

%%— Mise en page ---------------------------------------------------------------
\usepackage[margin=2.5cm]{geometry}
\usepackage{fancyhdr}
\usepackage{booktabs}

%%— Listes, index, liens -------------------------------------------------------
\usepackage{enumitem}
\usepackage{makeidx}
\usepackage{hyperref}
\usepackage{xcolor}

\makeindex

%%— En-têtes et pieds de page --------------------------------------------------
\pagestyle{fancy}
\fancyhf{}
\fancyhead[LE]{\leftmark}
\fancyhead[RO]{\rightmark}
\fancyfoot[C]{\thepage}
\renewcommand{\headrulewidth}{0.4pt}

%%=============================================================================
%% Environnements de théorèmes (amsthm + tcolorbox)
%%=============================================================================
\theoremstyle{definition}
\newtheorem{definition}{Définition}[chapter]
\newtheorem{theoreme}{Théorème}[chapter]
\newtheorem{proposition}{Proposition}[chapter]
\newtheorem{lemme}{Lemme}[chapter]
\newtheorem{corollaire}{Corollaire}[chapter]
\newtheorem{remarque}{Remarque}[chapter]
\newtheorem{exemple}{Exemple}[chapter]
\newtheorem{exercice}{Exercice}[chapter]

%%— Habillage tcolorbox --------------------------------------------------------
\tcolorboxenvironment{definition}{
  colback=blue!5!white,
  colframe=blue!75!black,
  fonttitle=\bfseries,
  boxrule=0.8pt,
  arc=2pt,
  left=6pt, right=6pt, top=6pt, bottom=6pt,
  breakable
}

\tcolorboxenvironment{theoreme}{
  colback=red!5!white,
  colframe=red!75!black,
  fonttitle=\bfseries,
  boxrule=0.8pt,
  arc=2pt,
  left=6pt, right=6pt, top=6pt, bottom=6pt,
  breakable
}

\tcolorboxenvironment{proposition}{
  colback=green!5!white,
  colframe=green!50!black,
  fonttitle=\bfseries,
  boxrule=0.8pt,
  arc=2pt,
  left=6pt, right=6pt, top=6pt, bottom=6pt,
  breakable
}

%%=============================================================================
%% Commandes personnalisées
%%=============================================================================
\newcommand{\N}{\mathbb{N}}
\newcommand{\Z}{\mathbb{Z}}
\newcommand{\Q}{\mathbb{Q}}
\newcommand{\R}{\mathbb{R}}
\newcommand{\C}{\mathbb{C}}
\newcommand{\F}{\mathbb{F}}
\newcommand{\Gal}{\operatorname{Gal}}
\newcommand{\Aut}{\operatorname{Aut}}
\newcommand{\Hom}{\operatorname{Hom}}
\newcommand{\Ker}{\operatorname{Ker}}
\renewcommand{\Im}{\operatorname{Im}}
\newcommand{\Id}{\operatorname{Id}}
\newcommand{\car}{\operatorname{car}}
\newcommand{\irr}{\operatorname{irr}}
\newcommand{\Fix}{\operatorname{Fix}}
\newcommand{\pgcd}{\operatorname{pgcd}}
\renewcommand{\deg}{\operatorname{deg}}

%%— Numérotation des chapitres à partir de 8 ----------------------------------
\setcounter{chapter}{7}

%%=============================================================================
%% Page de titre
%%=============================================================================
%%=============================================================================
\begin{document}
\begin{titlepage}
\begin{tikzpicture}[remember picture, overlay]
  \fill[left color=blue!15!white, right color=blue!5!white]
    (current page.south west) rectangle (current page.north east);
  \fill[color=orange!70!yellow]
    ([yshift=-2cm]current page.north west) rectangle ([yshift=-2.3cm]current page.north east);
  \fill[color=blue!80!black, rounded corners=8pt]
    ([xshift=3cm, yshift=-4cm]current page.north west) rectangle
    ([xshift=-3cm, yshift=-10cm]current page.north east);
  \node[anchor=center, text=white, font=\Huge\bfseries]
    at ([yshift=-6cm]current page.north) {Alg\`ebre Abstraite II};
  \node[anchor=center, text=orange!80!yellow, font=\Large]
    at ([yshift=-7.5cm]current page.north) {Notes de Cours};
  \node[anchor=center, text=white!80, font=\large]
    at ([yshift=-8.8cm]current page.north) {Licence L3 --- 2025--2026};
  \node[anchor=center, text=blue!80!black, font=\Large\itshape]
    at ([yshift=-12cm]current page.north) {Ya\'e Ulrich Gaba};
  \draw[orange!70!yellow, line width=2pt]
    ([xshift=5cm, yshift=-13.5cm]current page.north west) --
    ([xshift=-5cm, yshift=-13.5cm]current page.north east);
  \node[anchor=center, text width=12cm, align=center, text=blue!60!black, font=\itshape]
    at ([yshift=-16cm]current page.north) {<<\,Depuis Galois, l'alg\`ebre n'est plus l'art de r\'esoudre des \'equations, mais l'\'etude des structures.\,>>\\[4pt]--- Nicolas Bourbaki};
  \node[anchor=south, text=gray, font=\small]
    at ([yshift=1.5cm]current page.south) {\today};
\end{tikzpicture}
\end{titlepage}
\tableofcontents

%%=============================================================================
%% Préface
%%=============================================================================
\chapter*{Préface}
\addcontentsline{toc}{chapter}{Préface}

Ce second volume d'algèbre abstraite est consacré à la théorie des corps et à la
théorie de Galois, l'une des plus belles réalisations des mathématiques du
dix-neuvième siècle.

L'histoire de cette théorie est indissociable de la figure tragique d'\emph{Évariste
Galois}\index{Galois, Évariste} (1811--1832). Né à Bourg-la-Reine, Galois
manifesta très tôt un génie mathématique exceptionnel, mais sa vie fut marquée
par une série d'infortunes : ses mémoires soumis à l'Académie des sciences furent
perdus ou incompris --- Cauchy égara le premier, Fourier mourut avant de lire le
second, et Poisson le jugea incompréhensible. Républicain ardent dans la France
de la Restauration, Galois fut emprisonné à deux reprises pour ses activités
politiques. Le 30 mai 1832, à l'âge de vingt ans, il fut mortellement blessé lors
d'un duel dont les circonstances exactes demeurent obscures. La veille de sa mort,
dans une lettre fiévreuse adressée à son ami Auguste Chevalier, il résuma
l'essentiel de ses découvertes mathématiques, griffonnant dans les marges :
\og{}Je n'ai pas le temps.\fg{}

Ce que Galois avait découvert, c'est un lien profond entre la structure des
extensions de corps et celle des groupes de permutations. En associant à toute
équation polynomiale un groupe --- le \emph{groupe de Galois}\index{groupe de
Galois} --- il fournit un critère précis pour déterminer si l'équation est
résoluble par radicaux. Ce faisant, il prouva l'impossibilité de résoudre par
radicaux l'équation générale de degré cinq ou plus, un problème qui avait occupé
les plus grands algébristes pendant des siècles.

L'héritage de Galois dépasse largement la question de la résolubilité des
équations. Sa théorie est devenue un outil fondamental en théorie algébrique des
nombres, en géométrie algébrique et dans bien d'autres domaines. Le concept même
de \emph{groupe}, aujourd'hui central dans toutes les mathématiques, trouve l'une
de ses origines dans ses travaux.

Ce cours suppose acquis le contenu du premier volume (groupes, anneaux, idéaux,
anneaux de polynômes). Nous développerons la théorie des extensions de corps
(chapitres~8--10), puis la correspondance de Galois proprement dite
(chapitres~11--14).

\bigskip
\hfill\emph{L'auteur}

%%=============================================================================
%% Notations
%%=============================================================================
\chapter*{Notations}
\addcontentsline{toc}{chapter}{Notations}

\begin{tabular}{@{}ll@{}}
\toprule
\textbf{Notation} & \textbf{Signification} \\
\midrule
$\N, \Z, \Q, \R, \C$ & Entiers naturels, entiers relatifs, rationnels, réels, complexes \\
$\F_q$ & Corps fini à $q$ éléments \\
$K^*$, $K^\times$ & Groupe multiplicatif d'un corps $K$ \\
$L/K$ & Extension de corps ($K \subset L$) \\
$[L:K]$ & Degré de l'extension $L/K$ \\
$K(\alpha)$, $K(\alpha_1,\dots,\alpha_n)$ & Corps engendré par les $\alpha_i$ sur $K$ \\
$K[\alpha]$ & Sous-anneau de $L$ engendré par $\alpha$ sur $K$ \\
$\irr(\alpha,K)$ & Polynôme minimal de $\alpha$ sur $K$ \\
$\car(K)$ & Caractéristique du corps $K$ \\
$\overline{K}$ & Clôture algébrique de $K$ \\
$\Gal(L/K)$ & Groupe de Galois de l'extension $L/K$ \\
$\Aut(L)$ & Groupe des automorphismes de $L$ \\
$\Aut_K(L)$ & $K$-automorphismes de $L$ (synonyme de $\Gal(L/K)$) \\
$\Hom_K(L,M)$ & $K$-morphismes de corps de $L$ dans $M$ \\
$K[X]$ & Anneau des polynômes à coefficients dans $K$ \\
$K(X)$ & Corps des fractions rationnelles sur $K$ \\
$\deg(f)$ & Degré du polynôme $f$ \\
$f'$ & Dérivée formelle du polynôme $f$ \\
$\Phi_n$ & $n$-ième polynôme cyclotomique \\
$\mu_n$ & Groupe des racines $n$-ièmes de l'unité \\
$\operatorname{Fix}(H)$ & Corps fixe du groupe $H$ \\
$\operatorname{pgcd}$ & Plus grand commun diviseur \\
\bottomrule
\end{tabular}

%%#############################################################################
%%
%% CHAPITRE 8 : Corps — Extensions Algébriques et Transcendantes
%%
%%#############################################################################
\chapter{Corps --- Extensions algébriques et transcendantes}
\label{chap:extensions}

\section*{Introduction}
\addcontentsline{toc}{section}{Introduction}

Depuis les travaux des algébristes italiens du seizième siècle (Cardano, Tartaglia,
Ferrari) et la démonstration par Abel et Galois de l'impossibilité de résoudre par
radicaux l'équation générale de degré cinq, la question de la résolution des
équations polynomiales a guidé le développement de l'algèbre. Pour comprendre les
racines d'un polynôme $f \in K[X]$, il est naturel de considérer un corps plus
grand $L \supset K$ contenant ces racines : c'est la notion d'\emph{extension de
corps}\index{extension de corps}.

Dans ce chapitre, nous posons les fondements de la théorie des extensions de corps.
Nous étudierons la distinction entre éléments algébriques et transcendants, nous
montrerons que l'adjonction d'un élément algébrique $\alpha$ à un corps $K$
produit un corps isomorphe au quotient $K[X]/(\irr(\alpha,K))$, et nous
établirons les propriétés fondamentales des extensions finies et algébriques. Nous
aborderons également la caractéristique d'un corps et l'endomorphisme de Frobenius,
outils essentiels pour l'étude des corps finis.

%%=============================================================================
\section{Corps et sous-corps}
\label{sec:corps-sous-corps}

\begin{definition}[Corps]\label{def:corps}
\index{corps}
Un \emph{corps} est un anneau commutatif unitaire $(K, +, \cdot)$ dans lequel
tout élément non nul est inversible, c'est-à-dire $K^* = K \setminus \{0\}$
est un groupe pour la multiplication. On demande en outre $1 \neq 0$, ce qui
exclut l'anneau nul.
\end{definition}

\begin{definition}[Sous-corps]\label{def:sous-corps}
\index{sous-corps}
Soit $K$ un corps. Un sous-ensemble $F \subset K$ est un \emph{sous-corps} de
$K$ si $F$ est lui-même un corps pour les lois induites, c'est-à-dire :
\begin{enumerate}[label=(\roman*)]
  \item $1 \in F$ ;
  \item pour tous $a, b \in F$, on a $a - b \in F$ ;
  \item pour tous $a, b \in F$ avec $b \neq 0$, on a $ab^{-1} \in F$.
\end{enumerate}
\end{definition}

\begin{definition}[Extension de corps]\label{def:extension}
\index{extension de corps}
Une \emph{extension de corps} est un couple $(K, L)$ de corps tel que $K$ est
un sous-corps de $L$. On note $L/K$ et on dit que $L$ est une extension de $K$,
ou que $K$ est le \emph{corps de base}\index{corps de base}.
\end{definition}

\begin{remarque}[Structure de $K$-espace vectoriel]
\label{rem:espace-vectoriel}
Si $L/K$ est une extension de corps, alors $L$ est naturellement un $K$-espace
vectoriel pour l'addition de $L$ et la multiplication par les scalaires de $K$.
Cette observation, simple mais fondamentale, permet d'utiliser toute la puissance
de l'algèbre linéaire dans l'étude des extensions.
\end{remarque}

%%=============================================================================
\section{Degré d'une extension}
\label{sec:degre}

\begin{definition}[Degré]\label{def:degre}
\index{degré d'une extension}
Soit $L/K$ une extension de corps. Le \emph{degré} de l'extension, noté $[L:K]$,
est la dimension de $L$ en tant que $K$-espace vectoriel :
\[
  [L:K] = \dim_K(L).
\]
Si $[L:K] < \infty$, on dit que l'extension est \emph{finie}\index{extension
finie}. Sinon, elle est \emph{infinie}.
\end{definition}

\begin{exemple}[Extensions classiques]\label{ex:extensions-classiques}
\begin{enumerate}[label=(\alph*)]
  \item $[\C:\R] = 2$, car $(1,i)$ est une $\R$-base de $\C$.
  \item $[\R:\Q] = \infty$ (par un argument de cardinalité).
  \item $[\F_{p^n}:\F_p] = n$ pour tout premier $p$ et tout $n \geq 1$.
\end{enumerate}
\end{exemple}

\begin{theoreme}[Loi de multiplicativité des degrés]\label{thm:tour}
\index{tour (loi de la)}
Soit $K \subset L \subset M$ une tour d'extensions de corps. Alors
\[
  [M:K] = [M:L]\,[L:K].
\]
En particulier, $[M:K]$ est fini si et seulement si $[M:L]$ et $[L:K]$ sont
tous deux finis.
\end{theoreme}

\begin{proof}
Soit $(e_i)_{i \in I}$ une $K$-base de $L$ et $(f_j)_{j \in J}$ une $L$-base
de $M$. Montrons que la famille $(e_i f_j)_{(i,j) \in I \times J}$ est une
$K$-base de $M$.

\textbf{Famille génératrice.} Soit $m \in M$. Puisque $(f_j)_{j \in J}$ est
une $L$-base de $M$, il existe des éléments $\lambda_j \in L$ (presque tous
nuls) tels que $m = \sum_{j \in J} \lambda_j f_j$. Pour chaque $j$, puisque
$(e_i)_{i \in I}$ est une $K$-base de $L$, il existe des éléments
$\mu_{ij} \in K$ (presque tous nuls) tels que $\lambda_j = \sum_{i \in I}
\mu_{ij} e_i$. Par conséquent,
\[
  m = \sum_{j \in J} \left(\sum_{i \in I} \mu_{ij} e_i\right) f_j
    = \sum_{(i,j) \in I \times J} \mu_{ij}\, e_i f_j.
\]
La famille $(e_i f_j)$ est donc génératrice.

\textbf{Famille libre.} Supposons que
$\sum_{(i,j) \in I \times J} \mu_{ij}\, e_i f_j = 0$ avec $\mu_{ij} \in K$.
On réécrit cette relation :
\[
  \sum_{j \in J} \left(\sum_{i \in I} \mu_{ij} e_i\right) f_j = 0.
\]
Pour chaque $j$, posons $\lambda_j = \sum_{i \in I} \mu_{ij} e_i \in L$.
Puisque $(f_j)_{j \in J}$ est une $L$-base de $M$ (donc libre sur $L$), on a
$\lambda_j = 0$ pour tout $j \in J$. Mais alors, puisque $(e_i)_{i \in I}$ est
une $K$-base de $L$ (donc libre sur $K$), on obtient $\mu_{ij} = 0$ pour tous
$i, j$.

La famille $(e_i f_j)_{(i,j) \in I \times J}$ est donc une $K$-base de $M$, et
$[M:K] = |I \times J| = |I| \cdot |J| = [L:K]\,[M:L]$.
\end{proof}

\begin{center}
\begin{tikzcd}[row sep=1.5cm, column sep=0pt]
  M \arrow[d, dash, "{[M:L]}"'] \\
  L \arrow[d, dash, "{[L:K]}"'] \\
  K
\end{tikzcd}
\quad
\text{et}\quad $[M:K] = [M:L]\cdot [L:K]$.
\end{center}

\begin{corollaire}[Divisibilité du degré]\label{cor:divisibilite-degre}
Si $K \subset L \subset M$ et $[M:K]$ est fini, alors $[L:K]$ divise $[M:K]$.
\end{corollaire}

\begin{proof}
Par le théorème~\ref{thm:tour}, $[M:K] = [M:L][L:K]$, donc $[L:K]$ divise
$[M:K]$.
\end{proof}

%%=============================================================================
\section{Éléments algébriques et polynôme minimal}
\label{sec:elements-algebriques}

\begin{definition}[Élément algébrique, élément transcendant]\label{def:algebrique}
\index{élément algébrique}\index{élément transcendant}
Soit $L/K$ une extension de corps et $\alpha \in L$.
\begin{enumerate}[label=(\roman*)]
  \item On dit que $\alpha$ est \emph{algébrique sur $K$} s'il existe un
    polynôme non nul $f \in K[X]$ tel que $f(\alpha) = 0$.
  \item Sinon, $\alpha$ est dit \emph{transcendant sur $K$}.
\end{enumerate}
\end{definition}

\begin{exemple}[Éléments algébriques et transcendants]\label{ex:alg-trans}
\begin{enumerate}[label=(\alph*)]
  \item $\sqrt{2}$ est algébrique sur $\Q$ (racine de $X^2 - 2$).
  \item $i \in \C$ est algébrique sur $\R$ (racine de $X^2 + 1$).
  \item $\pi$ et $e$ sont transcendants sur $\Q$ (théorèmes de Lindemann et
    Hermite--Lindemann).
\end{enumerate}
\end{exemple}

\begin{definition}[Polynôme minimal]\label{def:polymin}
\index{polynôme minimal}
Soit $L/K$ une extension et $\alpha \in L$ algébrique sur $K$. Le \emph{polynôme
minimal} de $\alpha$ sur $K$, noté $\irr(\alpha, K)$, est le polynôme unitaire
de plus petit degré dans $K[X]$ qui annule $\alpha$.
\end{definition}

\begin{theoreme}[Existence, unicité et irréductibilité du polynôme minimal]
\label{thm:polymin}
\index{polynôme minimal!existence et unicité}
Soit $L/K$ une extension et $\alpha \in L$ algébrique sur $K$. Alors :
\begin{enumerate}[label=(\roman*)]
  \item Le polynôme minimal $\mu = \irr(\alpha, K)$ existe et est unique.
  \item $\mu$ est irréductible dans $K[X]$.
  \item Pour tout $f \in K[X]$, on a $f(\alpha) = 0$ si et seulement si
    $\mu \mid f$ dans $K[X]$.
\end{enumerate}
\end{theoreme}

\begin{proof}
Considérons le morphisme d'évaluation
$\varphi_\alpha \colon K[X] \to L$ défini par $f \mapsto f(\alpha)$. Le noyau
$\Ker(\varphi_\alpha) = \{f \in K[X] : f(\alpha) = 0\}$ est un idéal de $K[X]$.
Puisque $\alpha$ est algébrique sur $K$, cet idéal est non nul. Or $K[X]$ est un
anneau principal, donc $\Ker(\varphi_\alpha) = (\mu)$ pour un unique polynôme
unitaire $\mu \in K[X]$.

\textbf{(iii)} Par construction, $f(\alpha) = 0$ si et seulement si
$f \in (\mu)$, c'est-à-dire $\mu \mid f$.

\textbf{(i)} Le polynôme $\mu$ est unitaire et de degré minimal parmi les
polynômes non nuls qui annulent $\alpha$ (car tout élément de $(\mu)$ a un degré
$\geq \deg(\mu)$ ou est nul). L'unicité résulte de la condition d'unitarité et de
la minimalité du degré : si $\mu'$ était un autre tel polynôme, alors $\mu \mid
\mu'$ et $\mu' \mid \mu$, donc $\mu = \mu'$ (les deux étant unitaires).

\textbf{(ii)} Supposons $\mu = gh$ avec $g, h \in K[X]$ et $\deg(g), \deg(h) \geq
1$. Alors $0 = \mu(\alpha) = g(\alpha)h(\alpha)$. Puisque $L$ est un corps (donc
intègre), on a $g(\alpha) = 0$ ou $h(\alpha) = 0$. Donc $\mu \mid g$ ou $\mu \mid
h$ par (iii). Mais $\deg(g) < \deg(\mu)$ et $\deg(h) < \deg(\mu)$, contradiction.
Donc $\mu$ est irréductible.
\end{proof}

%%=============================================================================
\section{L'isomorphisme fondamental $K(\alpha) \cong K[X]/(\irr(\alpha,K))$}
\label{sec:iso-fondamental}

\begin{definition}[Corps engendré]\label{def:corps-engendre}
\index{corps engendré}
Soit $L/K$ une extension et $\alpha \in L$. On note $K(\alpha)$ le plus petit
sous-corps de $L$ contenant $K$ et $\alpha$. On note $K[\alpha]$ le plus petit
sous-anneau de $L$ contenant $K$ et $\alpha$, c'est-à-dire
\[
  K[\alpha] = \{f(\alpha) : f \in K[X]\}.
\]
\end{definition}

\begin{theoreme}[Isomorphisme fondamental]\label{thm:iso-fondamental}
\index{isomorphisme fondamental}
Soit $L/K$ une extension et $\alpha \in L$ algébrique sur $K$, de polynôme
minimal $\mu = \irr(\alpha, K)$ avec $n = \deg(\mu)$. Alors :
\begin{enumerate}[label=(\roman*)]
  \item $K[\alpha] = K(\alpha)$.
  \item L'application $\overline{\varphi} \colon K[X]/(\mu) \to K(\alpha)$
    définie par $\overline{f} \mapsto f(\alpha)$ est un isomorphisme de corps.
  \item $(1, \alpha, \alpha^2, \dots, \alpha^{n-1})$ est une $K$-base de
    $K(\alpha)$.
  \item $[K(\alpha):K] = n = \deg(\irr(\alpha, K))$.
\end{enumerate}
\end{theoreme}

\begin{proof}
Considérons le morphisme d'évaluation
$\varphi_\alpha \colon K[X] \to L$, $f \mapsto f(\alpha)$. Son image est
$K[\alpha]$ et son noyau est $(\mu)$ par le théorème~\ref{thm:polymin}.

Par le premier théorème d'isomorphisme,
\[
  K[X]/(\mu) \;\cong\; K[\alpha].
\]
Puisque $\mu$ est irréductible dans $K[X]$ (anneau principal), l'idéal $(\mu)$
est maximal, donc $K[X]/(\mu)$ est un corps. Il s'ensuit que $K[\alpha]$ est un
corps.

\textbf{(i)} Puisque $K[\alpha]$ est un corps contenant $K$ et $\alpha$, et que
$K(\alpha)$ est le plus petit tel corps, on a $K(\alpha) \subset K[\alpha]$.
L'inclusion réciproque $K[\alpha] \subset K(\alpha)$ est évidente, d'où
$K[\alpha] = K(\alpha)$.

\textbf{(ii)} L'isomorphisme $\overline{\varphi} \colon K[X]/(\mu) \xrightarrow{\sim} K(\alpha)$
découle du premier théorème d'isomorphisme.

\textbf{(iii)} Les classes $\overline{1}, \overline{X}, \dots, \overline{X^{n-1}}$
forment une $K$-base de $K[X]/(\mu)$ (car tout élément de $K[X]/(\mu)$ admet un
unique représentant de degré $< n = \deg(\mu)$). Via l'isomorphisme
$\overline{\varphi}$, la famille $1, \alpha, \dots, \alpha^{n-1}$ est une $K$-base
de $K(\alpha)$.

\textbf{(iv)} Il résulte de (iii) que $[K(\alpha):K] = n = \deg(\mu)$.
\end{proof}

\begin{exemple}[$\Q(\sqrt{2})$]\label{ex:Q-sqrt2}
\index{Q(sqrt2)@$\Q(\sqrt{2})$}
Le polynôme $X^2 - 2$ est irréductible dans $\Q[X]$ (critère d'Eisenstein avec
$p = 2$, ou absence de racine rationnelle). Donc $\irr(\sqrt{2}, \Q) = X^2 - 2$
et
\[
  \Q(\sqrt{2}) = \{a + b\sqrt{2} : a, b \in \Q\} \cong \Q[X]/(X^2 - 2),
  \qquad [\Q(\sqrt{2}):\Q] = 2.
\]
\end{exemple}

\begin{exemple}[{$\Q(\sqrt[3]{2})$}]\label{ex:Q-cbrt2}
\index{Q(cbrt2)@$\Q(\sqrt[3]{2})$}
Le polynôme $X^3 - 2$ est irréductible dans $\Q[X]$ (Eisenstein avec $p = 2$).
Donc $\irr(\sqrt[3]{2}, \Q) = X^3 - 2$ et
\[
  \Q(\sqrt[3]{2}) = \{a + b\sqrt[3]{2} + c\sqrt[3]{4} : a, b, c \in \Q\},
  \qquad [\Q(\sqrt[3]{2}):\Q] = 3.
\]
\end{exemple}

\begin{exemple}[$\Q(i)$]\label{ex:Q-i}
\index{Q(i)@$\Q(i)$}
Le polynôme $X^2 + 1$ est irréductible dans $\Q[X]$ (pas de racine rationnelle).
Donc
\[
  \Q(i) = \{a + bi : a, b \in \Q\} \cong \Q[X]/(X^2+1), \qquad [\Q(i):\Q] = 2.
\]
\end{exemple}

%%=============================================================================
\section{Extensions algébriques et extensions finies}
\label{sec:ext-algebriques}

\begin{definition}[Extension algébrique]\label{def:ext-algebrique}
\index{extension algébrique}
Une extension $L/K$ est dite \emph{algébrique} si tout élément de $L$ est
algébrique sur $K$.
\end{definition}

\begin{theoreme}[Finie implique algébrique]\label{thm:finie-implique-algebrique}
Toute extension finie est algébrique.
\end{theoreme}

\begin{proof}
Soit $L/K$ une extension finie avec $[L:K] = n$. Soit $\alpha \in L$. Les
$n+1$ éléments $1, \alpha, \alpha^2, \dots, \alpha^n$ appartiennent au $K$-espace
vectoriel $L$ de dimension $n$, donc ils sont $K$-linéairement dépendants. Il
existe ainsi $a_0, a_1, \dots, a_n \in K$, non tous nuls, tels que
$a_0 + a_1\alpha + \cdots + a_n\alpha^n = 0$. Le polynôme
$f = a_0 + a_1 X + \cdots + a_n X^n \in K[X]$ est non nul et $f(\alpha) = 0$,
donc $\alpha$ est algébrique sur $K$.
\end{proof}

\begin{remarque}[La réciproque est fausse]
\label{rem:algebrique-pas-finie}
La clôture algébrique $\overline{\Q}$ de $\Q$ est une extension algébrique de
$\Q$, mais elle n'est pas finie (car $[\Q(\sqrt[n]{2}):\Q] = n$ pour tout $n$).
\end{remarque}

\begin{theoreme}[La composée d'extensions algébriques est algébrique]
\label{thm:algebrique-algebrique}
Soient $K \subset L \subset M$ des corps. Si $L/K$ est algébrique et $M/L$
est algébrique, alors $M/K$ est algébrique.
\end{theoreme}

\begin{proof}
Soit $\alpha \in M$. Puisque $M/L$ est algébrique, $\alpha$ est algébrique sur
$L$, c'est-à-dire il existe un polynôme
$f = b_0 + b_1 X + \cdots + b_m X^m \in L[X]$ avec $b_m \neq 0$ et
$f(\alpha) = 0$. Posons $K' = K(b_0, b_1, \dots, b_m)$.

Puisque $L/K$ est algébrique, chaque $b_j$ est algébrique sur $K$. Par
adjonctions successives :
\[
  [K':K] = [K(b_0, \dots, b_m):K] \leq \prod_{j=0}^{m} \deg(\irr(b_j, K)) < \infty.
\]
Plus précisément, $K' = K(b_0)(b_1)\cdots(b_m)$ et chaque adjonction est de
degré fini par le théorème~\ref{thm:iso-fondamental}, puis la loi de
multiplicativité des degrés (théorème~\ref{thm:tour}) donne $[K':K] < \infty$.

Or $\alpha$ est racine de $f \in K'[X]$, donc $\alpha$ est algébrique sur $K'$
et $[K'(\alpha):K'] \leq m$. Par la loi de multiplicativité :
\[
  [K'(\alpha):K] = [K'(\alpha):K'][K':K] < \infty.
\]
Puisque $K \subset K'(\alpha)$ est une extension finie, elle est algébrique
(théorème~\ref{thm:finie-implique-algebrique}), et en particulier $\alpha$ est
algébrique sur $K$.
\end{proof}

\begin{corollaire}[L'ensemble des éléments algébriques est un corps]
\label{cor:elements-algebriques-corps}
Soit $L/K$ une extension. L'ensemble $\overline{K}^L = \{\alpha \in L :
\alpha \text{ est algébrique sur } K\}$ est un sous-corps de $L$.
\end{corollaire}

\begin{proof}
Soient $\alpha, \beta \in \overline{K}^L$. L'extension $K(\alpha,\beta)/K$ est
finie (car $[K(\alpha,\beta):K] \leq [K(\alpha):K][K(\beta):K] < \infty$), donc
algébrique. En particulier, $\alpha + \beta$, $\alpha\beta$ et, si
$\beta \neq 0$, $\alpha/\beta$ sont algébriques sur $K$.
\end{proof}

%%=============================================================================
\section{Extensions et éléments transcendants}
\label{sec:transcendant}

\begin{definition}[Extension transcendante]\label{def:ext-transcendante}
\index{extension transcendante}
Une extension $L/K$ est \emph{transcendante} si elle n'est pas algébrique,
c'est-à-dire s'il existe au moins un élément de $L$ qui est transcendant sur $K$.
\end{definition}

\begin{proposition}[Structure de $K(\alpha)$ pour $\alpha$ transcendant]
\label{prop:transcendant}
Soit $L/K$ une extension et $\alpha \in L$ transcendant sur $K$. Alors le
morphisme d'évaluation $\varphi_\alpha \colon K[X] \to L$, $f \mapsto f(\alpha)$,
est injectif, et il induit un isomorphisme $K(X) \cong K(\alpha)$ entre le corps
des fractions rationnelles et le sous-corps engendré par $\alpha$.
\end{proposition}

\begin{proof}
L'injectivité de $\varphi_\alpha$ résulte du fait que $\alpha$ est transcendant :
$\Ker(\varphi_\alpha) = \{f \in K[X] : f(\alpha) = 0\} = \{0\}$. Donc
$K[\alpha] \cong K[X]$. En passant aux corps de fractions (propriété universelle),
on obtient $K(\alpha) \cong K(X)$.
\end{proof}

\begin{remarque}[Degré de transcendance]
\label{rem:degre-transcendance}
\index{degré de transcendance}
Plus généralement, si $L/K$ est une extension, une \emph{base de transcendance}
de $L/K$ est un sous-ensemble $S \subset L$ maximal algébriquement indépendant
sur $K$. Le cardinal de $S$ (qui ne dépend pas du choix de $S$) est le
\emph{degré de transcendance} de $L/K$, noté $\operatorname{tr.deg}(L/K)$.
Par exemple, $\operatorname{tr.deg}(\R/\Q) = |\R|$ (indénombrable).
\end{remarque}

\begin{exemple}[L'extension $\F_p(t)/\F_p$]\label{ex:Fpt}
Soit $t$ une indéterminée. Le corps $\F_p(t)$ des fractions rationnelles sur
$\F_p$ est une extension transcendante de $\F_p$ : l'élément $t$ n'est racine
d'aucun polynôme non nul à coefficients dans $\F_p$. Le degré de transcendance
est~$1$.
\end{exemple}

%%=============================================================================
\section{Caractéristique et endomorphisme de Frobenius}
\label{sec:caracteristique}

\begin{definition}[Caractéristique]\label{def:caracteristique}
\index{caractéristique}
Soit $K$ un corps. La \emph{caractéristique} de $K$, notée $\car(K)$, est le
plus petit entier $p > 0$ tel que
$\underbrace{1 + 1 + \cdots + 1}_{p \text{ fois}} = 0$ dans $K$, s'il existe.
Sinon, $\car(K) = 0$.
\end{definition}

\begin{proposition}[La caractéristique est $0$ ou un nombre premier]
\label{prop:car-premier}
Pour tout corps $K$, $\car(K)$ est soit $0$, soit un nombre premier.
\end{proposition}

\begin{proof}
Si $\car(K) = n > 0$, supposons $n = ab$ avec $1 < a, b < n$. Alors dans $K$,
$(a \cdot 1)(b \cdot 1) = n \cdot 1 = 0$, où $a \cdot 1 = \underbrace{1 + \cdots + 1}_{a}$.
Puisque $K$ est intègre, $a \cdot 1 = 0$ ou $b \cdot 1 = 0$, ce qui contredit
la minimalité de $n$.
\end{proof}

\begin{proposition}[Sous-corps premier]\label{prop:sous-corps-premier}
\index{sous-corps premier}
Tout corps $K$ contient un unique plus petit sous-corps, appelé \emph{sous-corps
premier}, isomorphe à $\Q$ si $\car(K) = 0$, et à $\F_p$ si $\car(K) = p$.
\end{proposition}

\begin{proof}
L'image du morphisme d'anneaux canonique $\Z \to K$, $n \mapsto n \cdot 1_K$,
est un sous-anneau intègre de $K$. Si $\car(K) = 0$, ce morphisme est injectif,
son image est isomorphe à $\Z$ et le sous-corps premier est son corps de
fractions, isomorphe à $\Q$. Si $\car(K) = p$, le noyau est $p\Z$ et l'image
est isomorphe à $\Z/p\Z = \F_p$.
\end{proof}

\begin{theoreme}[Endomorphisme de Frobenius]\label{thm:frobenius}
\index{Frobenius (endomorphisme de)}
Soit $K$ un corps de caractéristique $p > 0$. L'application
\[
  \varphi \colon K \to K, \qquad x \mapsto x^p,
\]
est un morphisme de corps, appelé \emph{endomorphisme de Frobenius}.
\end{theoreme}

\begin{proof}
Il est clair que $\varphi(1) = 1$ et $\varphi(xy) = (xy)^p = x^p y^p =
\varphi(x)\varphi(y)$ (car $K$ est commutatif). Il reste à montrer que
$\varphi(x + y) = \varphi(x) + \varphi(y)$, c'est-à-dire
$(x+y)^p = x^p + y^p$.

Par la formule du binôme de Newton :
\[
  (x+y)^p = \sum_{k=0}^{p} \binom{p}{k} x^k y^{p-k}.
\]
Pour $1 \leq k \leq p-1$, le coefficient binomial
$\binom{p}{k} = \frac{p!}{k!(p-k)!}$ est un entier divisible par $p$ : en effet,
$p$ divise le numérateur $p!$ mais ne divise pas le dénominateur $k!(p-k)!$
(car $p$ est premier et $1 \leq k \leq p-1$). Puisque $\car(K) = p$, on a
$\binom{p}{k} \cdot 1_K = 0$ pour $1 \leq k \leq p-1$. Il reste
$(x+y)^p = x^p + y^p$.

Enfin, $\varphi$ est un morphisme de corps, donc injectif (un morphisme de
corps non nul est toujours injectif, car son noyau est un idéal de $K$ et les
seuls idéaux de $K$ sont $\{0\}$ et $K$).
\end{proof}

\begin{corollaire}[Itéré de Frobenius]\label{cor:frobenius-itere}
Pour tout $n \geq 1$, l'application $x \mapsto x^{p^n}$ est un morphisme de
corps de $K$ dans lui-même.
\end{corollaire}

\begin{proof}
C'est la composée $\varphi^n = \underbrace{\varphi \circ \cdots \circ \varphi}_{n}$,
composée de morphismes de corps.
\end{proof}

%%=============================================================================
\section{Corps finis}
\label{sec:corps-finis}

\begin{theoreme}[Existence et unicité des corps finis]\label{thm:corps-finis}
\index{corps fini!existence et unicité}
\begin{enumerate}[label=(\roman*)]
  \item Tout corps fini a $p^n$ éléments pour un certain nombre premier $p$ et
    un entier $n \geq 1$.
  \item Pour tout premier $p$ et tout $n \geq 1$, il existe un corps à $p^n$
    éléments.
  \item Deux corps finis de même cardinal sont isomorphes.
\end{enumerate}
On note $\F_{p^n}$\index{F_{p^n}@$\F_{p^n}$} l'unique corps à $p^n$ éléments
(à isomorphisme près).
\end{theoreme}

\begin{remarque}[Preuve reportée]
\label{rem:corps-finis-preuve}
La démonstration complète de l'existence et de l'unicité sera donnée au
chapitre~\ref{chap:cloture}, où nous montrerons que $\F_{p^n}$ est le corps
de décomposition de $X^{p^n} - X$ sur $\F_p$.
\end{remarque}

\begin{exemple}[Le corps $\F_4$]\label{ex:F4}
Le polynôme $X^2 + X + 1$ est irréductible sur $\F_2$ (il n'a pas de racine
dans $\F_2$). Soit $\alpha$ une racine dans une extension. Alors
\[
  \F_4 = \F_2(\alpha) = \{0, 1, \alpha, \alpha + 1\}
\]
avec $\alpha^2 = \alpha + 1$ (puisque $\car = 2$).
\end{exemple}

%%=============================================================================
\section{Exercices}
\label{sec:exercices-chap8}

\begin{exercice}[Degré de $\Q(\sqrt{2}, \sqrt{3})$]\label{exo:Q-sqrt2-sqrt3}
Montrer que $[\Q(\sqrt{2}, \sqrt{3}):\Q] = 4$ et déterminer une $\Q$-base de
$\Q(\sqrt{2}, \sqrt{3})$.
\emph{Indication :} montrer que $\sqrt{3} \notin \Q(\sqrt{2})$.
\end{exercice}

\begin{exercice}[Degré et divisibilité]\label{exo:degre-divisibilite}
Soit $L/K$ une extension de degré premier $p$. Montrer qu'il n'existe aucun
corps $F$ tel que $K \subsetneq F \subsetneq L$.
\end{exercice}

\begin{exercice}[Polynôme minimal de $\sqrt{2}+\sqrt{3}$]\label{exo:sqrt2+sqrt3}
\begin{enumerate}[label=(\alph*)]
  \item Montrer que $\Q(\sqrt{2}+\sqrt{3}) = \Q(\sqrt{2},\sqrt{3})$.
  \item En déduire $\irr(\sqrt{2}+\sqrt{3},\Q)$ et son degré.
\end{enumerate}
\end{exercice}

\begin{exercice}[{Polynôme minimal de $\sqrt[3]{2}$}]\label{exo:cbrt2-polymin}
Montrer que $X^3 - 2$ est irréductible sur $\Q$ et calculer $[\Q(\sqrt[3]{2}):\Q]$.
\end{exercice}

\begin{exercice}[Extension de degré 2]\label{exo:ext-degre2}
Soit $L/K$ une extension de degré $2$. Montrer que $L = K(\alpha)$ pour tout
$\alpha \in L \setminus K$.
\end{exercice}

\begin{exercice}[Sous-corps de $\C$]\label{exo:sous-corps-C}
Le corps $\Q(\sqrt[4]{2})$ est-il un sous-corps de $\R$ ? De $\C$ ?
Calculer $[\Q(\sqrt[4]{2}):\Q]$ et $[\Q(\sqrt[4]{2}):\Q(\sqrt{2})]$.
\end{exercice}

\begin{exercice}[Frobenius et corps parfait]\label{exo:frobenius}
Soit $K$ un corps fini de caractéristique $p$.
\begin{enumerate}[label=(\alph*)]
  \item Montrer que le Frobenius $\varphi \colon x \mapsto x^p$ est un
    automorphisme de $K$.
  \item En déduire que tout élément de $K$ est une puissance $p$-ième.
\end{enumerate}
\end{exercice}

\begin{exercice}[Corps de caractéristique $p$ non parfait]\label{exo:non-parfait}
Montrer que le Frobenius $\varphi \colon \F_p(t) \to \F_p(t)$, $x \mapsto x^p$,
n'est pas surjectif. En déduire que $\F_p(t)$ n'est pas parfait.
\end{exercice}

\begin{exercice}[$\F_4$ et $\F_8$]\label{exo:F4-F8}
\begin{enumerate}[label=(\alph*)]
  \item Construire explicitement $\F_8$ comme $\F_2[X]/(f)$ pour un polynôme
    irréductible $f$ de degré $3$ sur $\F_2$.
  \item Montrer que $\F_4$ n'est pas un sous-corps de $\F_8$.
\end{enumerate}
\end{exercice}

\begin{exercice}[Irréductibilité et extensions]\label{exo:irred-ext}
Soit $f \in K[X]$ irréductible de degré $n$. Soit $L/K$ une extension de degré
$m$ avec $\operatorname{pgcd}(n,m) = 1$. Montrer que $f$ est irréductible dans
$L[X]$.
\end{exercice}

\begin{exercice}[Tour d'extensions]\label{exo:tour}
Soit $K = \Q$, $L = \Q(\sqrt{2})$, $M = \Q(\sqrt[4]{2})$.
\begin{enumerate}[label=(\alph*)]
  \item Vérifier que $[M:L] = 2$, $[L:K] = 2$, $[M:K] = 4$.
  \item Trouver une $K$-base de $M$.
  \item Calculer $\irr(\sqrt[4]{2}, K)$ et $\irr(\sqrt[4]{2}, L)$.
\end{enumerate}
\end{exercice}

\bigskip
\begin{center}
\fbox{\parbox{0.85\textwidth}{
\textbf{Résumé du chapitre~\thechapter.}
\begin{itemize}
  \item Une extension $L/K$ fait de $L$ un $K$-espace vectoriel ; sa dimension
    est le degré $[L:K]$.
  \item La loi de multiplicativité des degrés : $[M:K]=[M:L][L:K]$.
  \item Un élément algébrique $\alpha$ a un polynôme minimal $\irr(\alpha,K)$
    irréductible et unitaire, et $K(\alpha) \cong K[X]/(\irr(\alpha,K))$.
  \item Finie $\Rightarrow$ algébrique ; composée d'algébriques = algébrique.
  \item En caractéristique $p$, le Frobenius $x \mapsto x^p$ est un
    endomorphisme de corps.
\end{itemize}
}}
\end{center}

%%#############################################################################
%%
%% CHAPITRE 9 : Corps de Rupture et Corps de Décomposition
%%
%%#############################################################################
\chapter{Corps de rupture et corps de décomposition}
\label{chap:decomposition}

\section*{Introduction}
\addcontentsline{toc}{section}{Introduction}

Au chapitre précédent, nous avons vu qu'un polynôme irréductible $f \in K[X]$
n'a pas nécessairement de racine dans $K$, mais que l'on peut construire une
extension $K(\alpha) \cong K[X]/(f)$ dans laquelle $f$ possède au moins une
racine. Cette construction soulève naturellement deux questions :
\begin{enumerate}[label=(\arabic*)]
  \item Peut-on construire une extension dans laquelle $f$ possède \emph{toutes}
    ses racines, c'est-à-dire se scinde en facteurs linéaires ?
  \item Une telle extension est-elle essentiellement unique ?
\end{enumerate}
Le présent chapitre répond affirmativement à ces deux questions. Nous introduisons
les notions de \emph{corps de rupture} et de \emph{corps de décomposition}, nous
démontrons leur existence et leur unicité (à isomorphisme près), puis nous abordons
la notion fondamentale de \emph{séparabilité}, qui jouera un rôle central dans la
théorie de Galois.

%%=============================================================================
\section{Corps de rupture}
\label{sec:corps-rupture}

\begin{definition}[Corps de rupture]\label{def:corps-rupture}
\index{corps de rupture}
Soit $K$ un corps et $f \in K[X]$ un polynôme irréductible. Un \emph{corps de
rupture} de $f$ sur $K$ est une extension $L/K$ de la forme $L = K(\alpha)$ où
$\alpha$ est une racine de $f$ dans $L$.
\end{definition}

\begin{proposition}[Existence du corps de rupture]\label{prop:existence-rupture}
Pour tout polynôme irréductible $f \in K[X]$, le corps $L = K[X]/(f)$ est un
corps de rupture de $f$ sur $K$.
\end{proposition}

\begin{proof}
Puisque $f$ est irréductible dans l'anneau principal $K[X]$, l'idéal $(f)$ est
maximal et $L = K[X]/(f)$ est un corps. L'application canonique $K \to L$,
$a \mapsto \overline{a}$, est un morphisme de corps (donc injectif), ce qui
permet d'identifier $K$ à un sous-corps de $L$. L'élément $\alpha =
\overline{X} \in L$ vérifie $f(\alpha) = \overline{f(X)} = \overline{0} = 0$
et $L = K(\alpha)$.
\end{proof}

%%=============================================================================
\section{Corps de décomposition}
\label{sec:corps-decomposition}

\begin{definition}[Corps de décomposition]\label{def:corps-decomposition}
\index{corps de décomposition}
Soit $K$ un corps et $f \in K[X]$ un polynôme non constant de degré $n$.
Un \emph{corps de décomposition}\footnote{On dit aussi \emph{corps de
scindage}.} de $f$ sur $K$ est une extension $L/K$ telle que :
\begin{enumerate}[label=(\roman*)]
  \item $f$ se scinde dans $L[X]$ : il existe $\alpha_1, \dots, \alpha_n \in L$
    et $a \in K^*$ tels que $f = a(X - \alpha_1)\cdots(X - \alpha_n)$ ;
  \item $L = K(\alpha_1, \dots, \alpha_n)$ (minimalité).
\end{enumerate}
\end{definition}

\begin{theoreme}[Existence du corps de décomposition]\label{thm:existence-decomp}
\index{corps de décomposition!existence}
Pour tout corps $K$ et tout polynôme non constant $f \in K[X]$, il existe un
corps de décomposition de $f$ sur $K$.
\end{theoreme}

\begin{proof}
On procède par récurrence sur $n = \deg(f)$.

\textbf{Cas de base : $n = 1$.} Si $\deg(f) = 1$, alors $f = a(X - \alpha)$
avec $\alpha \in K$, et $K$ lui-même est un corps de décomposition de $f$.

\textbf{Étape de récurrence.} Supposons $n \geq 2$ et le résultat vrai pour
tout polynôme de degré $< n$ sur tout corps. Soit $g$ un facteur irréductible
de $f$ dans $K[X]$, de degré $d \geq 1$. Par la proposition~\ref{prop:existence-rupture},
il existe une extension $K_1 = K[X]/(g)$ dans laquelle $g$ possède une racine
$\alpha_1$. On a donc $f = (X - \alpha_1) f_1$ dans $K_1[X]$, avec
$\deg(f_1) = n - 1$.

Par hypothèse de récurrence, il existe un corps de décomposition $L$ de $f_1$
sur $K_1$. Soient $\alpha_2, \dots, \alpha_n \in L$ les racines de $f_1$ dans $L$
(comptées avec multiplicité). Alors
$f = a(X - \alpha_1)(X - \alpha_2)\cdots(X - \alpha_n)$ dans $L[X]$ et
\[
  L = K_1(\alpha_2, \dots, \alpha_n) = K(\alpha_1, \alpha_2, \dots, \alpha_n)
\]
est un corps de décomposition de $f$ sur $K$.
\end{proof}

%%=============================================================================
\section{Extension de morphismes de corps}
\label{sec:extension-morphismes}

Pour démontrer l'unicité du corps de décomposition, nous avons besoin d'un lemme
fondamental sur l'extension des morphismes de corps.

\begin{lemme}[Extension d'un isomorphisme]\label{lem:extension-iso}
\index{extension d'isomorphisme}
Soit $\sigma \colon K \xrightarrow{\sim} K'$ un isomorphisme de corps. Soit
$f \in K[X]$ irréductible, et soit $f^\sigma \in K'[X]$ le polynôme obtenu en
appliquant $\sigma$ aux coefficients de $f$. Soit $\alpha$ une racine de $f$
dans une extension de $K$ et $\beta$ une racine de $f^\sigma$ dans une extension
de $K'$. Alors il existe un unique isomorphisme
$\tau \colon K(\alpha) \xrightarrow{\sim} K'(\beta)$ tel que
$\tau|_K = \sigma$ et $\tau(\alpha) = \beta$.
\end{lemme}

\begin{proof}
Posons $n = \deg(f)$. Par le théorème~\ref{thm:iso-fondamental}, tout élément
de $K(\alpha)$ s'écrit de manière unique sous la forme
$a_0 + a_1 \alpha + \cdots + a_{n-1}\alpha^{n-1}$ avec $a_i \in K$, et de même
tout élément de $K'(\beta)$ s'écrit de manière unique sous la forme
$b_0 + b_1 \beta + \cdots + b_{n-1}\beta^{n-1}$ avec $b_i \in K'$.

Définissons $\tau \colon K(\alpha) \to K'(\beta)$ par
\[
  \tau\!\left(\sum_{i=0}^{n-1} a_i \alpha^i\right) = \sum_{i=0}^{n-1} \sigma(a_i)\, \beta^i.
\]
Vérifions que $\tau$ est un morphisme de corps. La $K$-linéarité (et donc
l'additivité) est claire par construction.

Pour la multiplicativité, considérons deux éléments $u = \sum a_i \alpha^i$ et
$v = \sum b_j \alpha^j$. Leur produit dans $K(\alpha)$ est obtenu en calculant
$\left(\sum a_i X^i\right)\left(\sum b_j X^j\right)$ puis en réduisant modulo
$f$ (qui est le polynôme minimal de $\alpha$ sur $K$). De même, $\tau(u)\tau(v)$
est obtenu en calculant
$\left(\sum \sigma(a_i) X^i\right)\left(\sum \sigma(b_j) X^j\right)$ et en
réduisant modulo $f^\sigma$ (polynôme minimal de $\beta$ sur $K'$). Puisque
$\sigma$ est un isomorphisme de corps, il commute avec toutes les opérations
arithmétiques sur les coefficients, et la réduction modulo $f$ chez $K$ correspond
exactement à la réduction modulo $f^\sigma$ chez $K'$. Donc
$\tau(uv) = \tau(u)\tau(v)$.

Par construction, $\tau|_K = \sigma$ et $\tau(\alpha) = \beta$.

L'injectivité est automatique (morphisme de corps). La surjectivité résulte du
fait que $\tau$ envoie la $K$-base $(1, \alpha, \dots, \alpha^{n-1})$ sur la
$K'$-base $(1, \beta, \dots, \beta^{n-1})$.

L'unicité est claire : tout morphisme $\tau' \colon K(\alpha) \to K'(\beta)$
avec $\tau'|_K = \sigma$ et $\tau'(\alpha) = \beta$ vérifie
$\tau'(\sum a_i \alpha^i) = \sum \sigma(a_i)\beta^i = \tau(\sum a_i \alpha^i)$.
\end{proof}

\begin{center}
\begin{tikzcd}[column sep=2cm, row sep=1.5cm]
  K(\alpha) \arrow[r, "\tau", dashed, "\sim"'] & K'(\beta) \\
  K \arrow[u, hook] \arrow[r, "\sigma"', "\sim"] & K' \arrow[u, hook]
\end{tikzcd}
\end{center}

%%=============================================================================
\section{Unicité du corps de décomposition}
\label{sec:unicite-decomp}

\begin{theoreme}[Unicité du corps de décomposition]\label{thm:unicite-decomp}
\index{corps de décomposition!unicité}
Soit $\sigma \colon K \xrightarrow{\sim} K'$ un isomorphisme de corps. Soit
$f \in K[X]$ non constant et $f^\sigma \in K'[X]$ le polynôme correspondant.
Soit $L$ un corps de décomposition de $f$ sur $K$ et $L'$ un corps de
décomposition de $f^\sigma$ sur $K'$. Alors il existe un isomorphisme
$\tau \colon L \xrightarrow{\sim} L'$ qui prolonge $\sigma$ (i.e.\
$\tau|_K = \sigma$).
\end{theoreme}

\begin{proof}
On procède par récurrence sur $n = \deg(f)$.

\textbf{Cas de base : $n = 1$.} Alors $L = K$ et $L' = K'$, et $\tau = \sigma$
convient.

\textbf{Étape de récurrence.} Supposons $n \geq 2$ et le résultat vrai pour
tout polynôme de degré $< n$. Soit $g$ un facteur irréductible de $f$ dans $K[X]$
et soit $g^\sigma$ le facteur irréductible correspondant de $f^\sigma$ dans
$K'[X]$. Soit $\alpha \in L$ une racine de $g$ et $\beta \in L'$ une racine
de $g^\sigma$.

Par le lemme~\ref{lem:extension-iso}, il existe un isomorphisme
$\sigma_1 \colon K(\alpha) \xrightarrow{\sim} K'(\beta)$ prolongeant $\sigma$
avec $\sigma_1(\alpha) = \beta$.

Dans $K(\alpha)[X]$, on a $f = (X - \alpha)\, f_1$ et $L$ est un corps de
décomposition de $f_1$ sur $K(\alpha)$. De même, dans $K'(\beta)[X]$, on a
$f^\sigma = (X - \beta)\, f_1^{\sigma_1}$ et $L'$ est un corps de décomposition
de $f_1^{\sigma_1}$ sur $K'(\beta)$. Puisque $\deg(f_1) = n - 1 < n$,
l'hypothèse de récurrence fournit un isomorphisme
$\tau \colon L \xrightarrow{\sim} L'$ prolongeant $\sigma_1$, et donc
prolongeant $\sigma$.
\end{proof}

\begin{corollaire}[Unicité à isomorphisme près]\label{cor:unicite-decomp}
Pour tout corps $K$ et tout $f \in K[X]$ non constant, le corps de décomposition
de $f$ sur $K$ est unique à $K$-isomorphisme près.
\end{corollaire}

\begin{proof}
Appliquer le théorème~\ref{thm:unicite-decomp} avec $\sigma = \Id_K$.
\end{proof}

%%=============================================================================
\section{Exemples de corps de décomposition}
\label{sec:exemples-decomp}

\begin{exemple}[Corps de décomposition de $X^2 + 1$ sur $\R$]\label{ex:decomp-X2+1}
Le polynôme $X^2 + 1$ n'a pas de racine dans $\R$. En posant $\C = \R[X]/(X^2+1)$
et $i = \overline{X}$, on obtient $X^2 + 1 = (X - i)(X + i)$ dans $\C[X]$ et
$\C = \R(i, -i) = \R(i)$. Ainsi $\C$ est le corps de décomposition de $X^2+1$
sur $\R$.
\end{exemple}

\begin{exemple}[Corps de décomposition de $X^3 - 2$ sur $\Q$]\label{ex:decomp-X3-2}
\index{X3-2@$X^3 - 2$}
Les trois racines de $X^3 - 2$ dans $\C$ sont $\sqrt[3]{2}$,
$\sqrt[3]{2}\,\omega$ et $\sqrt[3]{2}\,\omega^2$, où
$\omega = e^{2i\pi/3} = \frac{-1+i\sqrt{3}}{2}$.

Le corps de décomposition est $L = \Q(\sqrt[3]{2}, \omega)$. Calculons son degré.
\begin{itemize}
  \item $[\Q(\sqrt[3]{2}):\Q] = 3$ (car $X^3-2$ est irréductible sur $\Q$).
  \item $\omega$ est racine de $X^2 + X + 1$ (irréductible sur $\Q$, donc aussi
    sur $\Q(\sqrt[3]{2})$ puisque $3 \nmid 2$). Ainsi
    $[L:\Q(\sqrt[3]{2})] = 2$.
  \item Par la loi multiplicative : $[L:\Q] = 3 \times 2 = 6$.
\end{itemize}
\end{exemple}

\begin{center}
\begin{tikzcd}[row sep=1.2cm, column sep=0pt]
  & \Q(\sqrt[3]{2},\omega) \arrow[dl, dash, "2"'] \arrow[dr, dash, "3"] & \\
  \Q(\sqrt[3]{2}) \arrow[dr, dash, "3"'] & & \Q(\omega) \arrow[dl, dash, "2"] \\
  & \Q &
\end{tikzcd}
\end{center}

\begin{exemple}[Corps cyclotomique]\label{ex:corps-cyclotomique}
\index{corps cyclotomique}\index{polynôme cyclotomique}
Le corps de décomposition de $X^n - 1$ sur $\Q$ est le \emph{corps cyclotomique}
$\Q(\zeta_n)$, où $\zeta_n = e^{2i\pi/n}$ est une racine primitive $n$-ième de
l'unité. En effet, les racines de $X^n - 1$ sont $1, \zeta_n, \zeta_n^2, \dots,
\zeta_n^{n-1}$, toutes puissances de $\zeta_n$, donc
$\Q(1, \zeta_n, \dots, \zeta_n^{n-1}) = \Q(\zeta_n)$.

Le polynôme minimal de $\zeta_n$ sur $\Q$ est le $n$-ième polynôme cyclotomique
$\Phi_n \in \Z[X]$, de degré $\varphi(n)$ (indicatrice d'Euler). Ainsi
$[\Q(\zeta_n):\Q] = \varphi(n)$.
\end{exemple}

%%=============================================================================
\section{Racines multiples et dérivée formelle}
\label{sec:racines-multiples}

\begin{definition}[Multiplicité d'une racine]\label{def:multiplicite}
\index{multiplicité}
Soit $K$ un corps, $f \in K[X]$ non nul et $\alpha$ une racine de $f$ dans une
extension de $K$. La \emph{multiplicité} de $\alpha$ est le plus grand entier
$m \geq 1$ tel que $(X - \alpha)^m \mid f$ dans $\overline{K}[X]$. Si $m = 1$,
$\alpha$ est une \emph{racine simple}\index{racine simple} ; si $m \geq 2$,
$\alpha$ est une \emph{racine multiple}\index{racine multiple}.
\end{definition}

\begin{definition}[Dérivée formelle]\label{def:derivee-formelle}
\index{dérivée formelle}
Soit $f = \sum_{k=0}^n a_k X^k \in K[X]$. La \emph{dérivée formelle} de $f$ est
\[
  f' = \sum_{k=1}^{n} k\, a_k X^{k-1} \in K[X].
\]
\end{definition}

\begin{proposition}[Critère de racine multiple]\label{prop:critere-racine-multiple}
Soit $f \in K[X]$ non constant. Un élément $\alpha$ (dans une extension de $K$)
est racine multiple de $f$ si et seulement si $f(\alpha) = 0$ et $f'(\alpha) = 0$.
Plus généralement, $f$ a une racine multiple (dans une clôture algébrique) si et
seulement si $\gcd(f, f') \neq 1$.
\end{proposition}

\begin{proof}
Écrivons $f = (X - \alpha)^m g$ avec $g(\alpha) \neq 0$. Alors
$f' = m(X-\alpha)^{m-1}g + (X-\alpha)^m g'$. Si $m \geq 2$,
$(X-\alpha) \mid f'$, donc $f'(\alpha) = 0$. Si $m = 1$,
$f' = g + (X-\alpha)g'$, donc $f'(\alpha) = g(\alpha) \neq 0$.
\end{proof}

%%=============================================================================
\section{Polynômes séparables et extensions séparables}
\label{sec:separable}

\begin{definition}[Polynôme séparable]\label{def:poly-separable}
\index{polynôme séparable}
Un polynôme $f \in K[X]$ irréductible est \emph{séparable} s'il n'a pas de
racine multiple dans une clôture algébrique de $K$, c'est-à-dire si
$\gcd(f, f') = 1$.

Un polynôme non constant quelconque est séparable si tous ses facteurs
irréductibles sont séparables.
\end{definition}

\begin{definition}[Élément séparable, extension séparable]\label{def:ext-separable}
\index{extension séparable}\index{élément séparable}
Soit $L/K$ une extension algébrique.
\begin{enumerate}[label=(\roman*)]
  \item Un élément $\alpha \in L$ est \emph{séparable sur $K$} si son polynôme
    minimal $\irr(\alpha, K)$ est séparable.
  \item L'extension $L/K$ est \emph{séparable} si tout élément de $L$ est
    séparable sur $K$.
\end{enumerate}
\end{definition}

\begin{theoreme}[En caractéristique $0$, toute extension algébrique est séparable]
\label{thm:car0-separable}
Si $\car(K) = 0$, alors tout polynôme irréductible $f \in K[X]$ est séparable.
En conséquence, toute extension algébrique de $K$ est séparable.
\end{theoreme}

\begin{proof}
Soit $f \in K[X]$ irréductible de degré $n \geq 1$. Alors
$f' \neq 0$ car $\deg(f') = n - 1 \geq 0$ (le coefficient dominant de $f'$ est
$n \cdot a_n$ où $a_n$ est le coefficient dominant de $f$, et $n \cdot a_n \neq 0$
en caractéristique $0$). Puisque $f$ est irréductible et $0 < \deg(f') < \deg(f)$,
le pgcd $\gcd(f, f')$ ne peut être un multiple non trivial de $f$. Comme $f$ est
irréductible, ses seuls diviseurs sont les constantes et les associés de $f$.
Puisque $\deg(f') < \deg(f)$, on a $\gcd(f, f') = 1$, donc $f$ est séparable.
\end{proof}

\begin{remarque}[Inséparabilité en caractéristique $p$]
\label{rem:inseparable}
En caractéristique $p > 0$, un polynôme irréductible $f$ peut être inséparable.
Cela se produit exactement lorsque $f' = 0$, c'est-à-dire lorsque $f$ est un
polynôme en $X^p$. Par exemple, sur $\F_p(t)$, le polynôme $X^p - t$ est
irréductible et inséparable.
\end{remarque}

%%=============================================================================
\section{Corps parfaits}
\label{sec:parfaits}

\begin{definition}[Corps parfait]\label{def:corps-parfait}
\index{corps parfait}
Un corps $K$ est \emph{parfait} si toute extension algébrique de $K$ est séparable.
De manière équivalente, $K$ est parfait si tout polynôme irréductible dans $K[X]$
est séparable.
\end{definition}

\begin{proposition}[Caractérisation des corps parfaits]\label{prop:parfait-car}
\begin{enumerate}[label=(\roman*)]
  \item Tout corps de caractéristique $0$ est parfait.
  \item Un corps $K$ de caractéristique $p > 0$ est parfait si et seulement si
    le Frobenius $\varphi \colon K \to K$, $x \mapsto x^p$, est surjectif
    (c'est-à-dire tout élément de $K$ est une puissance $p$-ième).
\end{enumerate}
\end{proposition}

\begin{proof}
\textbf{(i)} C'est le théorème~\ref{thm:car0-separable}.

\textbf{(ii)} Supposons $\car(K) = p$. Si $f \in K[X]$ est irréductible et
inséparable, alors $f' = 0$, ce qui signifie que $f = \sum_{i} a_i X^{ip}$,
c'est-à-dire $f$ est un polynôme en $X^p$. Si le Frobenius est surjectif,
chaque $a_i = b_i^p$ pour un $b_i \in K$, et alors
$f = \left(\sum_i b_i X^i\right)^p$ (en caractéristique $p$), ce qui contredit
l'irréductibilité de $f$. Donc tout irréductible est séparable.

Réciproquement, si $a \in K$ n'est pas une puissance $p$-ième, le polynôme
$X^p - a$ est irréductible dans $K[X]$ (admis) et sa dérivée est $pX^{p-1} = 0$,
donc il est inséparable.
\end{proof}

\begin{corollaire}[Les corps finis sont parfaits]\label{cor:finis-parfaits}
Tout corps fini est parfait.
\end{corollaire}

\begin{proof}
Le Frobenius $\varphi \colon K \to K$, $x \mapsto x^p$, est un morphisme
de corps injectif. Si $K$ est fini, toute injection $K \to K$ est une bijection,
donc $\varphi$ est surjectif.
\end{proof}

%%=============================================================================
\section{Exercices}
\label{sec:exercices-chap9}

\begin{exercice}[Corps de décomposition de $X^4 - 2$]\label{exo:decomp-X4-2}
Déterminer le corps de décomposition de $X^4 - 2$ sur $\Q$ et calculer son
degré sur $\Q$.
\end{exercice}

\begin{exercice}[Corps de décomposition de $X^3 - 1$]\label{exo:decomp-X3-1}
Déterminer le corps de décomposition de $X^3 - 1$ sur $\Q$, sur $\F_2$ et
sur $\F_7$.
\end{exercice}

\begin{exercice}[Extension de morphismes]\label{exo:extension-morphismes}
Soit $\sigma \colon \Q \to \Q$ l'identité et $f = X^3 - 2$. Déterminer tous
les prolongements de $\sigma$ en un morphisme $\Q(\sqrt[3]{2}) \to \C$.
Combien y en a-t-il ?
\end{exercice}

\begin{exercice}[Corps de décomposition sur $\F_p$]\label{exo:decomp-Fp}
Soit $p$ premier et $f = X^{p^2} - X \in \F_p[X]$.
\begin{enumerate}[label=(\alph*)]
  \item Montrer que $f$ se scinde dans $\F_{p^2}$.
  \item Montrer que $f$ n'a pas de racine multiple.
  \item En déduire que $\F_{p^2}$ est le corps de décomposition de $f$ sur
    $\F_p$.
\end{enumerate}
\end{exercice}

\begin{exercice}[Séparabilité]\label{exo:separabilite}
\begin{enumerate}[label=(\alph*)]
  \item Montrer que $X^4 + 1$ est séparable sur $\Q$.
  \item Montrer que $X^p - t \in \F_p(t)[X]$ est irréductible et inséparable.
\end{enumerate}
\end{exercice}

\begin{exercice}[Polynôme cyclotomique $\Phi_p$]\label{exo:cyclotomique}
Soit $p$ premier. Montrer que $\Phi_p(X) = X^{p-1} + X^{p-2} + \cdots + X + 1$
est irréductible sur $\Q$.
\emph{Indication :} considérer $\Phi_p(X+1)$ et appliquer le critère d'Eisenstein.
\end{exercice}

\begin{exercice}[Nombre de racines]\label{exo:nombre-racines}
Soit $f \in K[X]$ de degré $n$. Montrer que $f$ a au plus $n$ racines dans $K$
(comptées avec multiplicité).
\end{exercice}

\begin{exercice}[Dérivée formelle]\label{exo:derivee-formelle}
Montrer que les règles habituelles de dérivation sont valables pour la dérivée
formelle : $(f+g)' = f' + g'$, $(fg)' = f'g + fg'$, $(f^n)' = nf^{n-1}f'$.
\end{exercice}

\begin{exercice}[Corps de rupture vs.\ corps de décomposition]\label{exo:rupture-vs-decomp}
Soit $f = X^3 - 2 \in \Q[X]$.
\begin{enumerate}[label=(\alph*)]
  \item Montrer que $\Q(\sqrt[3]{2})$ est un corps de rupture de $f$ mais pas
    un corps de décomposition.
  \item Montrer que $\Q(\sqrt[3]{2}, \omega)$ est le corps de décomposition.
\end{enumerate}
\end{exercice}

\begin{exercice}[Unicité et automorphismes]\label{exo:unicite-auto}
Soit $L$ le corps de décomposition de $X^3 - 2$ sur $\Q$. Combien de
$\Q$-automorphismes $L \to L$ existe-t-il ? Les décrire explicitement.
\end{exercice}

\begin{exercice}[Séparabilité et extension finie]\label{exo:sep-ext-finie}
Soit $K$ un corps parfait et $L/K$ une extension finie. Montrer que $L/K$ est
séparable.
\end{exercice}

\begin{exercice}[Corps parfait et Frobenius]\label{exo:parfait-frobenius}
Montrer que $\F_p$ est parfait et que $\F_p(t)$ n'est pas parfait.
\end{exercice}

\bigskip
\begin{center}
\fbox{\parbox{0.85\textwidth}{
\textbf{Résumé du chapitre~\thechapter.}
\begin{itemize}
  \item Le corps de rupture d'un irréductible $f$ est $K[X]/(f)$.
  \item Le corps de décomposition de $f$ existe (récurrence sur le degré) et
    est unique à $K$-isomorphisme près.
  \item Le lemme d'extension des isomorphismes est l'outil clé de l'unicité.
  \item Un polynôme irréductible est séparable ssi $\gcd(f,f') = 1$.
  \item En caractéristique $0$, tout irréductible est séparable. Un corps est
    parfait ssi le Frobenius est surjectif (en char.\ $p$).
\end{itemize}
}}
\end{center}

%%#############################################################################
%%
%% CHAPITRE 10 : Clôture Algébrique
%%
%%#############################################################################
\chapter{Clôture algébrique}
\label{chap:cloture}

\section*{Introduction}
\addcontentsline{toc}{section}{Introduction}

Nous avons vu au chapitre précédent que, pour tout polynôme $f \in K[X]$, on peut
construire un corps de décomposition dans lequel $f$ possède toutes ses racines.
Mais peut-on construire une extension \emph{universelle} de $K$ dans laquelle
\emph{tout} polynôme non constant de $K[X]$ possède une racine ? Une telle
extension, appelée \emph{clôture algébrique}, est l'objet de ce chapitre.

Nous montrerons l'existence (qui repose sur le lemme de Zorn) et l'unicité (à
isomorphisme près) de la clôture algébrique, puis nous utiliserons cet outil pour
développer la théorie des corps finis de manière complète.

%%=============================================================================
\section{Corps algébriquement clos}
\label{sec:algebriquement-clos}

\begin{definition}[Corps algébriquement clos]\label{def:algebriquement-clos}
\index{corps algébriquement clos}
Un corps $K$ est \emph{algébriquement clos} si tout polynôme non constant de
$K[X]$ a (au moins) une racine dans $K$. De manière équivalente, les seuls
polynômes irréductibles de $K[X]$ sont les polynômes de degré~$1$.
\end{definition}

\begin{proposition}[Caractérisations équivalentes]\label{prop:car-alg-clos}
Soit $K$ un corps. Les conditions suivantes sont équivalentes :
\begin{enumerate}[label=(\roman*)]
  \item $K$ est algébriquement clos.
  \item Tout $f \in K[X]$ non constant se scinde dans $K[X]$.
  \item $K$ n'admet aucune extension algébrique propre.
\end{enumerate}
\end{proposition}

\begin{proof}
$\textbf{(i)} \Rightarrow \textbf{(ii)}$ :
Si $f \in K[X]$ est non constant, il a une racine $\alpha_1 \in K$, donc
$f = (X - \alpha_1) f_1$ avec $\deg(f_1) = \deg(f) - 1$. Par récurrence sur
le degré, $f$ se scinde dans $K[X]$.

$\textbf{(ii)} \Rightarrow \textbf{(iii)}$ :
Soit $L/K$ une extension algébrique et $\alpha \in L$. Son polynôme minimal
$\irr(\alpha, K)$ se scinde dans $K[X]$ par hypothèse, et en particulier
$\alpha \in K$ (car $\irr(\alpha, K)$ a une racine dans $K$, et étant
irréductible, il est de degré~$1$). Donc $L = K$.

$\textbf{(iii)} \Rightarrow \textbf{(i)}$ :
Soit $f \in K[X]$ irréductible. Le corps $K[X]/(f)$ est une extension algébrique
de $K$, de degré $\deg(f)$. Par hypothèse, $K[X]/(f) = K$, donc $\deg(f) = 1$.
\end{proof}

%%=============================================================================
\section{Clôture algébrique}
\label{sec:cloture}

\begin{definition}[Clôture algébrique]\label{def:cloture}
\index{clôture algébrique}
Soit $K$ un corps. Une \emph{clôture algébrique} de $K$ est un corps
$\overline{K}$ tel que :
\begin{enumerate}[label=(\roman*)]
  \item $\overline{K}/K$ est une extension algébrique ;
  \item $\overline{K}$ est algébriquement clos.
\end{enumerate}
\end{definition}

\begin{theoreme}[Existence de la clôture algébrique]\label{thm:existence-cloture}
\index{clôture algébrique!existence}
Tout corps $K$ admet une clôture algébrique.
\end{theoreme}

\begin{proof}[Esquisse de la preuve]
La preuve classique, due à Steinitz (1910), utilise le lemme de Zorn. Considérons
l'ensemble $\mathcal{E}$ des extensions algébriques de $K$ contenues dans un
univers ensembliste suffisamment grand (ou, de manière plus précise, on fixe un
ensemble $\Omega$ de cardinal $|\Omega| > |K[X]|$ et on considère les structures
de corps sur des sous-ensembles de $\Omega$ qui sont des extensions algébriques
de $K$).

On munit $\mathcal{E}$ de la relation d'ordre $L_1 \leq L_2$ si $L_1$ est un
sous-corps de $L_2$. Toute chaîne $(L_i)_{i \in I}$ de $\mathcal{E}$ admet un
majorant $\bigcup_{i \in I} L_i$ qui est encore une extension algébrique de $K$.
Par le lemme de Zorn, $\mathcal{E}$ admet un élément maximal $\overline{K}$.

Si $\overline{K}$ n'était pas algébriquement clos, il existerait un polynôme
irréductible $f \in \overline{K}[X]$ de degré $\geq 2$, et le corps
$\overline{K}[X]/(f)$ serait une extension algébrique stricte de $\overline{K}$,
donc une extension algébrique stricte de $K$ (la composée d'extensions algébriques
est algébrique, théorème~\ref{thm:algebrique-algebrique}), contredisant la
maximalité de $\overline{K}$.
\end{proof}

\begin{theoreme}[Unicité de la clôture algébrique]\label{thm:unicite-cloture}
\index{clôture algébrique!unicité}
Soient $\overline{K}$ et $\overline{K}'$ deux clôtures algébriques de $K$.
Alors il existe un $K$-isomorphisme
$\sigma \colon \overline{K} \xrightarrow{\sim} \overline{K}'$.
\end{theoreme}

\begin{proof}[Esquisse de la preuve]
Considérons l'ensemble des couples $(L, \sigma)$ où $L$ est un sous-corps de
$\overline{K}$ contenant $K$ et $\sigma \colon L \to \overline{K}'$ est un
$K$-morphisme de corps. Cet ensemble est non vide (il contient
$(K, \iota)$ où $\iota$ est l'inclusion $K \hookrightarrow \overline{K}'$) et
partiellement ordonné par $(L_1, \sigma_1) \leq (L_2, \sigma_2)$ si
$L_1 \subset L_2$ et $\sigma_2|_{L_1} = \sigma_1$.

Par le lemme de Zorn, il existe un élément maximal $(M, \tau)$. Montrons que
$M = \overline{K}$. Si $M \neq \overline{K}$, soit $\alpha \in \overline{K}
\setminus M$. Le polynôme $\irr(\alpha, M) \in M[X]$ est irréductible ; le
polynôme image $\irr(\alpha, M)^\tau \in \tau(M)[X]$ a une racine $\beta$ dans
$\overline{K}'$ (car $\overline{K}'$ est algébriquement clos). Par le
lemme~\ref{lem:extension-iso}, $\tau$ se prolonge en un morphisme
$\tau' \colon M(\alpha) \to \overline{K}'$ avec $\tau'(\alpha) = \beta$, ce qui
contredit la maximalité de $(M, \tau)$.

Donc $\tau \colon \overline{K} \to \overline{K}'$ est un $K$-morphisme de corps.
Montrons la surjectivité. L'image $\tau(\overline{K})$ est un sous-corps
algébriquement clos de $\overline{K}'$ (car isomorphe à $\overline{K}$) et
$\overline{K}'/\tau(\overline{K})$ est algébrique (car $\overline{K}'/K$ est
algébrique et $K \subset \tau(\overline{K})$). Par la
proposition~\ref{prop:car-alg-clos}(iii), $\tau(\overline{K}) = \overline{K}'$.
\end{proof}

%%=============================================================================
\section{Le théorème fondamental de l'algèbre}
\label{sec:thm-fondamental-algebre}

\begin{theoreme}[Théorème fondamental de l'algèbre]\label{thm:fta}
\index{théorème fondamental de l'algèbre}
Le corps $\C$ des nombres complexes est algébriquement clos.
\end{theoreme}

\begin{remarque}[Remarques historiques]
\label{rem:fta-histoire}
Ce résultat, pressenti par d'Alembert (1746), a été démontré rigoureusement pour
la première fois par Gauss dans sa thèse de doctorat (1799), bien que sa première
preuve comportât une lacune topologique. Gauss en donna ensuite trois autres
démonstrations. Malgré son nom, le théorème fondamental de l'algèbre est un
résultat d'analyse : toute preuve connue utilise une propriété de complétude de
$\R$ (ou de $\C$). Parmi les preuves classiques, citons :
\begin{enumerate}[label=(\alph*)]
  \item la preuve par le théorème de Liouville (analyse complexe) ;
  \item la preuve topologique utilisant le degré d'une application ;
  \item la preuve algébrique utilisant la théorie de Galois et le fait que $\R$
    est un corps réel clos.
\end{enumerate}
\end{remarque}

\begin{corollaire}[Clôture algébrique de $\R$]
\label{cor:cloture-R}
$\C$ est la clôture algébrique de $\R$, c'est-à-dire $\overline{\R} = \C$.
\end{corollaire}

\begin{proof}
$\C$ est algébriquement clos (théorème~\ref{thm:fta}) et $[\C:\R] = 2 < \infty$,
donc $\C/\R$ est algébrique (théorème~\ref{thm:finie-implique-algebrique}).
\end{proof}

%%=============================================================================
\section{La clôture algébrique de la clôture algébrique}
\label{sec:cloture-cloture}

\begin{proposition}[La clôture algébrique est sa propre clôture]
\label{prop:cloture-idempotent}
Soit $\overline{K}$ une clôture algébrique de $K$. Alors
$\overline{\overline{K}} = \overline{K}$.
\end{proposition}

\begin{proof}
Puisque $\overline{K}$ est algébriquement clos, il n'admet aucune extension
algébrique propre (proposition~\ref{prop:car-alg-clos}(iii)). Or toute clôture
algébrique de $\overline{K}$ est une extension algébrique de $\overline{K}$,
donc elle est égale à $\overline{K}$.
\end{proof}

%%=============================================================================
\section{Corps finis : existence, unicité et structure}
\label{sec:corps-finis-structure}

Nous sommes maintenant en mesure de démontrer complètement les résultats annoncés
au chapitre~\ref{chap:extensions} sur les corps finis.

\begin{theoreme}[Corps fini comme corps de décomposition]
\label{thm:corps-fini-decomp}
\index{corps fini!structure}
Soit $p$ un nombre premier et $n \geq 1$ un entier. Alors :
\begin{enumerate}[label=(\roman*)]
  \item L'ensemble des racines de $X^{p^n} - X$ dans $\overline{\F_p}$ forme
    un corps à $p^n$ éléments.
  \item Ce corps est le corps de décomposition de $X^{p^n} - X$ sur $\F_p$.
  \item Tout corps à $p^n$ éléments est isomorphe à ce corps.
\end{enumerate}
\end{theoreme}

\begin{proof}
\textbf{(i)} Posons $q = p^n$ et $f(X) = X^q - X$. Montrons d'abord que $f$
n'a pas de racine multiple. On a $f'(X) = qX^{q-1} - 1 = -1$ (car $q = p^n$
est un multiple de $p$ et $\car(\overline{\F_p}) = p$). Puisque $f' = -1 \neq 0$,
on a $\gcd(f, f') = 1$, donc $f$ a exactement $q = p^n$ racines distinctes dans
$\overline{\F_p}$.

Soit $S = \{\alpha \in \overline{\F_p} : \alpha^q = \alpha\}$ l'ensemble de ces
racines. Montrons que $S$ est un corps.
\begin{itemize}
  \item $0 \in S$ et $1 \in S$ (évidement).
  \item Si $\alpha, \beta \in S$, alors
    $(\alpha - \beta)^q = \alpha^q - \beta^q = \alpha - \beta$ (en
    caractéristique $p$, l'application $x \mapsto x^q = x^{p^n}$ est un
    morphisme de corps, corollaire~\ref{cor:frobenius-itere}), donc
    $\alpha - \beta \in S$.
  \item Si $\alpha, \beta \in S$, alors
    $(\alpha\beta)^q = \alpha^q \beta^q = \alpha\beta$, donc $\alpha\beta \in S$.
  \item Si $\alpha \in S$, $\alpha \neq 0$, alors
    $(\alpha^{-1})^q = (\alpha^q)^{-1} = \alpha^{-1}$, donc $\alpha^{-1} \in S$.
\end{itemize}
Ainsi $S$ est un sous-corps de $\overline{\F_p}$ à $q$ éléments.

\textbf{(ii)} Le polynôme $f = X^q - X$ se scinde dans $S[X]$ (car $S$ contient
toutes ses racines) et $S$ est engendré par ses racines sur $\F_p$ (car $S$ est
lui-même l'ensemble des racines). Donc $S$ est le corps de décomposition de $f$
sur $\F_p$.

\textbf{(iii)} Soit $K$ un corps à $q = p^n$ éléments. Alors $\car(K) = p$ et
$K^*$ est un groupe de cardinal $q - 1$, donc $\alpha^{q-1} = 1$ pour tout
$\alpha \in K^*$ (par le théorème de Lagrange). Ainsi $\alpha^q = \alpha$ pour
tout $\alpha \in K$ (y compris $\alpha = 0$), c'est-à-dire tout élément de $K$
est racine de $X^q - X$. Puisque $|K| = q = \deg(X^q - X)$ et qu'un polynôme
de degré $q$ a au plus $q$ racines, $K$ est exactement l'ensemble des racines de
$X^q - X$ dans toute extension, et donc $K$ est le corps de décomposition de
$X^q - X$ sur $\F_p$. Par l'unicité du corps de décomposition
(corollaire~\ref{cor:unicite-decomp}), $K \cong S$.
\end{proof}

\begin{theoreme}[Réseau des sous-corps d'un corps fini]\label{thm:sous-corps-fini}
\index{corps fini!sous-corps}
Soit $q = p^n$. Les sous-corps de $\F_{q}$ sont exactement les $\F_{p^d}$ pour
$d \mid n$, et $\F_{p^d} \subset \F_{p^e}$ si et seulement si $d \mid e$.
\end{theoreme}

\begin{proof}
\textbf{Condition nécessaire.} Soit $F$ un sous-corps de $\F_{p^n}$ avec
$|F| = p^d$. Alors $\F_{p^n}/F$ est une extension de degré
$[\F_{p^n}:F] = n/d$ (car $p^n = (p^d)^{n/d}$, et $[\F_{p^n}:\F_p] =
[\F_{p^n}:F][\F:\F_p]$, d'où $n = [\F_{p^n}:F] \cdot d$ et $[\F_{p^n}:F] =
n/d$). En particulier, $n/d$ est un entier, c'est-à-dire $d \mid n$.

\textbf{Condition suffisante.} Supposons $d \mid n$. Alors $p^d - 1 \mid p^n - 1$
(car $X^d - 1 \mid X^n - 1$ dans $\Z[X]$ quand $d \mid n$, évalué en $X = p$).
Donc $X^{p^d} - X \mid X^{p^n} - X$ dans $\F_p[X]$. En effet, si
$\alpha^{p^d} = \alpha$, alors
$\alpha^{p^n} = \alpha^{(p^d)^{n/d}} = \alpha$ (en itérant le Frobenius).
Ainsi toute racine de $X^{p^d} - X$ est racine de $X^{p^n} - X$, ce qui
signifie $\F_{p^d} \subset \F_{p^n}$.

\textbf{Inclusion.} Le même argument montre que $\F_{p^d} \subset \F_{p^e}$ si et
seulement si tout élément de $\F_{p^d}$ est racine de $X^{p^e} - X$, ce qui
équivaut à $p^d - 1 \mid p^e - 1$, ce qui équivaut à $d \mid e$.
\end{proof}

\begin{center}
\begin{tikzpicture}[scale=1, every node/.style={draw, circle, minimum size=1.2cm, inner sep=1pt, font=\small}]
  \node (1)  at (0,0)    {$\F_p$};
  \node (2)  at (-3,2)   {$\F_{p^2}$};
  \node (3)  at (-1,2)   {$\F_{p^3}$};
  \node (4)  at (1,2)    {$\F_{p^4}$};
  \node (6)  at (3,2)    {$\F_{p^6}$};
  \node (42) at (-2,4)   {$\F_{p^4}$};
  \node (62) at (0,4)    {$\F_{p^6}$};
  \node (12) at (2,4)    {$\F_{p^{12}}$};
  %% Redraw with correct lattice
\end{tikzpicture}
\end{center}

\vspace{-1cm}
Le réseau des sous-corps de $\F_{p^{12}}$ est le suivant. Les diviseurs de $12$
sont $1, 2, 3, 4, 6, 12$, et la relation d'inclusion est donnée par la
divisibilité :

\begin{center}
\begin{tikzcd}[row sep=1.2cm, column sep=0.8cm]
  & & \F_{p^{12}} \arrow[dll, dash] \arrow[dl, dash] \arrow[d, dash] \arrow[dr, dash] & \\
  \F_{p^{4}} \arrow[dr, dash] \arrow[drr, dash] &
  \F_{p^{6}} \arrow[d, dash] \arrow[dr, dash] &
  \F_{p^{12}} &
  {} \\
  & \F_{p^{2}} \arrow[dr, dash] &
  \F_{p^{3}} \arrow[d, dash] & \\
  & & \F_{p} &
\end{tikzcd}
\end{center}

Corrigeons et dessinons proprement le treillis :

\begin{center}
\begin{tikzcd}[row sep=1.5cm, column sep=1.5cm]
  & \F_{p^{12}} & \\
  \F_{p^{4}} \arrow[ur, dash] & \F_{p^{6}} \arrow[u, dash] & \\
  \F_{p^{2}} \arrow[u, dash] \arrow[ur, dash] & \F_{p^{3}} \arrow[u, dash] & \\
  & \F_{p} \arrow[ul, dash] \arrow[u, dash] &
\end{tikzcd}
\end{center}

%%=============================================================================
\section{Le groupe multiplicatif d'un corps fini est cyclique}
\label{sec:multiplicatif-cyclique}

\begin{theoreme}[Cyclicité du groupe multiplicatif]\label{thm:multiplicatif-cyclique}
\index{corps fini!groupe multiplicatif cyclique}
Soit $K$ un corps fini. Alors le groupe multiplicatif $K^*$ est cyclique.
\end{theoreme}

\begin{proof}
Soit $|K| = q$ et $K^* = K \setminus \{0\}$, de sorte que $|K^*| = q-1$. Nous
devons montrer que $K^*$ est cyclique, c'est-à-dire qu'il existe un élément
d'ordre $q - 1$ dans $K^*$.

Soit $n = q - 1$. Écrivons la factorisation en nombres premiers :
$n = p_1^{a_1} p_2^{a_2} \cdots p_r^{a_r}$.

\textbf{Étape 1 : existence d'un élément d'ordre $p_i^{a_i}$ pour chaque $i$.}
Fixons $i$ et posons $m = n / p_i^{a_i}$ et $d = p_i^{a_i}$. Le polynôme
$X^m - 1$ a au plus $m$ racines dans $K$. Puisque $m < n = |K^*|$, il existe un
élément $\beta \in K^*$ tel que $\beta^m \neq 1$. Posons
$\gamma = \beta^{m / p_i^{a_i - 1}} = \beta^{n / p_i}$ si $a_i \geq 1$.

Plus directement : l'élément $\beta^{n/p_i^{a_i}}$ a un ordre qui divise
$p_i^{a_i}$ (car $(\beta^{n/p_i^{a_i}})^{p_i^{a_i}} = \beta^n = 1$). Son ordre
est exactement $p_i^{a_i}$ si et seulement si $\beta^{n/p_i} \neq 1$.

Puisque $X^{n/p_i} - 1$ a au plus $n/p_i < n$ racines dans $K^*$, et
$|K^*| = n$, il existe $\beta \in K^*$ avec $\beta^{n/p_i} \neq 1$. Posons
$\gamma_i = \beta^{n / p_i^{a_i}}$. Alors l'ordre de $\gamma_i$ divise
$p_i^{a_i}$ (car $\gamma_i^{p_i^{a_i}} = \beta^n = 1$), et l'ordre de
$\gamma_i$ ne divise pas $p_i^{a_i - 1}$ (car
$\gamma_i^{p_i^{a_i - 1}} = \beta^{n/p_i} \neq 1$). Donc l'ordre de
$\gamma_i$ est exactement $p_i^{a_i}$.

\textbf{Étape 2 : construction d'un élément d'ordre $n$.}
Posons $g = \gamma_1 \gamma_2 \cdots \gamma_r$. Puisque les ordres
$p_1^{a_1}, \dots, p_r^{a_r}$ sont deux à deux premiers entre eux et que
$K^*$ est abélien, l'ordre de $g$ est
\[
  \operatorname{ord}(g) = \operatorname{ppcm}(p_1^{a_1}, \dots, p_r^{a_r}) =
  p_1^{a_1} \cdots p_r^{a_r} = n = q - 1.
\]
Pour justifier ce dernier point : dans un groupe abélien, si $\alpha$ est
d'ordre $a$ et $\beta$ est d'ordre $b$ avec $\gcd(a,b) = 1$, alors
$\alpha\beta$ est d'ordre $ab$. En effet, si $(\alpha\beta)^k = 1$, alors
$\alpha^k = \beta^{-k}$, et l'ordre de $\alpha^k$ divise à la fois $a$ et $b$,
donc est~$1$, d'où $\alpha^k = 1$ et $\beta^k = 1$, c'est-à-dire $a \mid k$ et
$b \mid k$, et finalement $ab \mid k$.

En appliquant ce résultat par récurrence aux $\gamma_i$ (dont les ordres sont
deux à deux premiers entre eux), on obtient
$\operatorname{ord}(g) = p_1^{a_1} \cdots p_r^{a_r} = n$.

Donc $g$ est un générateur de $K^*$, ce qui prouve que $K^*$ est cyclique.
\end{proof}

\begin{remarque}[Générateur primitif]
\label{rem:generateur-primitif}
Un générateur de $K^*$ est appelé \emph{élément primitif}\index{élément primitif}
ou \emph{racine primitive} de $K$. Par exemple, $2$ est un élément primitif de
$\F_5^*$ (car $2^1 = 2$, $2^2 = 4$, $2^3 = 3$, $2^4 = 1$), mais $4$ ne l'est
pas (car $4^2 = 1$).
\end{remarque}

\begin{corollaire}[Sous-groupes du groupe multiplicatif]\label{cor:sous-groupes-mult}
Soit $K$ un corps fini avec $|K| = q$. Pour tout diviseur $d$ de $q - 1$, il
existe un unique sous-groupe de $K^*$ d'ordre $d$, constitué des racines de
$X^d - 1$ dans $K$.
\end{corollaire}

\begin{proof}
Puisque $K^*$ est cyclique d'ordre $q - 1$, le résultat découle de la
classification des sous-groupes d'un groupe cyclique.
\end{proof}

%%=============================================================================
\section{Exercices}
\label{sec:exercices-chap10}

\begin{exercice}[Corps algébriquement clos]\label{exo:alg-clos}
Montrer qu'un corps fini n'est jamais algébriquement clos.
\emph{Indication :} considérer $f = 1 + \prod_{\alpha \in K} (X - \alpha)$.
\end{exercice}

\begin{exercice}[Clôture algébrique de $\F_p$]\label{exo:cloture-Fp}
Montrer que $\overline{\F_p} = \bigcup_{n \geq 1} \F_{p^{n!}}$.
\end{exercice}

\begin{exercice}[Sous-corps de $\F_{p^{12}}$]\label{exo:sous-corps-p12}
Lister tous les sous-corps de $\F_{p^{12}}$ et dessiner le diagramme de Hasse
de l'inclusion.
\end{exercice}

\begin{exercice}[Groupe multiplicatif de $\F_7$]\label{exo:mult-F7}
Déterminer un générateur de $\F_7^*$ et donner la table des puissances.
\end{exercice}

\begin{exercice}[Groupe multiplicatif de $\F_9$]\label{exo:mult-F9}
Construire $\F_9 = \F_3[X]/(X^2+1)$ et déterminer un générateur de $\F_9^*$.
\end{exercice}

\begin{exercice}[Irréductibles sur $\F_p$]\label{exo:irred-Fp}
\begin{enumerate}[label=(\alph*)]
  \item Montrer que $X^{p^n} - X = \prod_{d \mid n} \prod_f f$ dans $\F_p[X]$,
    le produit intérieur portant sur les polynômes irréductibles unitaires de
    degré $d$ dans $\F_p[X]$.
  \item En déduire que le nombre de polynômes irréductibles unitaires de degré
    $n$ dans $\F_p[X]$ est $\frac{1}{n}\sum_{d \mid n} \mu(n/d)\, p^d$, où
    $\mu$ est la fonction de Möbius.
\end{enumerate}
\end{exercice}

\begin{exercice}[Automorphismes de $\F_{p^n}$]\label{exo:aut-Fpn}
Montrer que $\Aut(\F_{p^n})$ est le groupe cyclique d'ordre $n$ engendré par
le Frobenius $\varphi \colon x \mapsto x^p$.
\end{exercice}

\begin{exercice}[Norme et trace sur les corps finis]\label{exo:norme-trace}
Soit $\F_{p^n}/\F_p$. Pour $\alpha \in \F_{p^n}$, on définit la trace
$\operatorname{Tr}(\alpha) = \sum_{k=0}^{n-1} \alpha^{p^k}$ et la norme
$\operatorname{N}(\alpha) = \prod_{k=0}^{n-1} \alpha^{p^k}$.
\begin{enumerate}[label=(\alph*)]
  \item Montrer que $\operatorname{Tr}(\alpha) \in \F_p$ et
    $\operatorname{N}(\alpha) \in \F_p$.
  \item Montrer que $\operatorname{Tr} \colon \F_{p^n} \to \F_p$ est
    $\F_p$-linéaire et surjective.
\end{enumerate}
\end{exercice}

\begin{exercice}[Existence d'éléments d'ordre donné]\label{exo:ordre-donne}
Soit $K$ un corps fini avec $|K^*| = q - 1$. Soit $d \mid q - 1$.
\begin{enumerate}[label=(\alph*)]
  \item Montrer qu'il existe exactement $\varphi(d)$ éléments d'ordre $d$ dans
    $K^*$.
  \item Retrouver la formule $\sum_{d \mid n} \varphi(d) = n$.
\end{enumerate}
\end{exercice}

\bigskip
\begin{center}
\fbox{\parbox{0.85\textwidth}{
\textbf{Résumé du chapitre~\thechapter.}
\begin{itemize}
  \item Un corps est algébriquement clos si tout polynôme non constant a une
    racine.
  \item Tout corps $K$ admet une clôture algébrique $\overline{K}$, unique à
    $K$-isomorphisme près.
  \item $\C$ est algébriquement clos (théorème fondamental de l'algèbre) et
    $\overline{\R} = \C$.
  \item $\F_{p^n}$ est le corps de décomposition de $X^{p^n} - X$ sur $\F_p$.
  \item Les sous-corps de $\F_{p^n}$ sont les $\F_{p^d}$ avec $d \mid n$.
  \item Le groupe multiplicatif d'un corps fini est cyclique.
\end{itemize}
}}
\end{center}
% ============================================================
%  Algèbre Abstraite II — Chapitres 11–13
%  Partie 2 (finale) : Extensions galoisiennes, Théorème
%  fondamental, Applications
% ============================================================

% ============================================================
%  CHAPITRE 11 : EXTENSIONS GALOISIENNES
% ============================================================
\chapter{Extensions galoisiennes}\label{chap:extensions-galoisiennes}
\index{extension!galoisienne}

\bigskip
\noindent
L'idée fondamentale d'Évariste Galois\index{Galois, Évariste}, exposée dans sa
célèbre lettre à Auguste Chevalier la veille de son duel fatal, est d'étudier
les racines d'un polynôme non pas individuellement mais à travers les
\emph{symétries} qui les permutent~: les automorphismes du corps engendré par
ces racines. Cette vision, radicalement nouvelle pour son époque, transforme un
problème d'algèbre en un problème de théorie des groupes. Le présent chapitre
développe cette machinerie~: nous définissons le groupe des automorphismes d'une
extension, isolons la classe des extensions pour lesquelles ce groupe encode
\emph{toute} l'information --- les extensions galoisiennes --- et calculons
explicitement de nombreux exemples.

% ------------------------------------------------------------
\section{Groupe d'automorphismes d'une extension}\label{sec:aut-extension}
\index{groupe!d'automorphismes}
% ------------------------------------------------------------

\begin{definition}[Automorphisme de corps]\label{def:aut-corps}
Soit $L/K$ une extension de corps\index{extension de corps}. Un
\emph{$K$-auto\-morphisme} de $L$ est un isomorphisme de corps
$\sigma \colon L \to L$ tel que $\sigma|_K = \mathrm{id}_K$, c'est-à-dire
$\sigma(a) = a$ pour tout $a \in K$. L'ensemble de tous les $K$-automorphismes
de $L$, muni de la composition, forme un groupe noté
\[
  \Aut(L/K) = \{\sigma \in \Aut(L) \mid \sigma|_K = \mathrm{id}_K\}.
\]
\end{definition}

\begin{remarque}[Premiers exemples]\label{rem:premiers-exemples-aut}
\begin{enumerate}[label=(\roman*)]
  \item Pour toute extension $L/K$, l'identité $\mathrm{id}_L$ appartient à
    $\Aut(L/K)$. On a toujours $|\Aut(L/K)| \geq 1$.
  \item Si $L/K$ est purement inséparable (et $L \neq K$), alors
    $\Aut(L/K) = \{\mathrm{id}_L\}$~: les automorphismes sont triviaux.
  \item Tout $\sigma \in \Aut(L/K)$ permute les racines de tout polynôme
    $f \in K[X]$ contenues dans $L$. En effet, si $f(\alpha) = 0$ et
    $f = \sum a_i X^i$ avec $a_i \in K$, alors
    $f(\sigma(\alpha)) = \sum a_i \sigma(\alpha)^i = \sigma(f(\alpha))
    = \sigma(0) = 0$.
\end{enumerate}
\end{remarque}

\begin{exemple}[$\Aut(\C/\R)$]\label{ex:aut-C-R}
\index{conjugaison complexe}
L'extension $\C/\R$ admet exactement deux $\R$-automorphismes~: l'identité et
la conjugaison complexe $z \mapsto \bar{z}$. En effet, tout
$\sigma \in \Aut(\C/\R)$ fixe $\R$ et doit envoyer $i$ sur une racine de
$X^2 + 1$ dans $\C$, soit $i$ ou $-i$. Comme $\sigma$ est entièrement déterminé
par $\sigma(i)$ (car $\C = \R(i)$), on obtient $|\Aut(\C/\R)| = 2$.
\end{exemple}

\begin{proposition}[Majoration du cardinal]\label{prop:majoration-aut}
Soit $L/K$ une extension finie de degré $n = [L:K]$. Alors
\[
  |\Aut(L/K)| \leq [L:K].
\]
\end{proposition}

\begin{proof}
Écrivons $L = K(\alpha_1, \ldots, \alpha_r)$ (ce qui est toujours possible,
$L/K$ étant de degré fini). Tout $\sigma \in \Aut(L/K)$ est entièrement
déterminé par les valeurs $\sigma(\alpha_1), \ldots, \sigma(\alpha_r)$.
Or $\sigma(\alpha_i)$ est une racine du polynôme minimal
$\irr(\alpha_i, K) \in K[X]$, dont le degré divise $[L:K]$. Par un argument
de comptage inductif sur les extensions successives
$K \subset K(\alpha_1) \subset K(\alpha_1,\alpha_2) \subset \cdots \subset L$,
le nombre de choix possibles est au plus
$[K(\alpha_1):K] \cdot [K(\alpha_1,\alpha_2):K(\alpha_1)] \cdots = [L:K]$.
\end{proof}

% ------------------------------------------------------------
\section{Corps fixe}\label{sec:corps-fixe}
\index{corps!fixe}
% ------------------------------------------------------------

\begin{definition}[Corps fixe]\label{def:corps-fixe}
Soit $L$ un corps et $H$ un sous-groupe de $\Aut(L)$. Le \emph{corps fixe}
de $H$ est
\[
  L^H = \Fix(H) = \{x \in L \mid \sigma(x) = x\ \text{pour tout }\
  \sigma \in H\}.
\]
C'est un sous-corps de $L$.
\end{definition}

\begin{proof}[Vérification]
Si $a, b \in L^H$ et $\sigma \in H$, alors
$\sigma(a + b) = \sigma(a) + \sigma(b) = a + b$,
$\sigma(ab) = \sigma(a)\sigma(b) = ab$,
$\sigma(-a) = -\sigma(a) = -a$, et pour $a \neq 0$,
$\sigma(a^{-1}) = \sigma(a)^{-1} = a^{-1}$. De plus $0, 1 \in L^H$.
Donc $L^H$ est un sous-corps.
\end{proof}

\begin{remarque}[Relation entre $K$ et $L^{\Aut(L/K)}$]\label{rem:corps-fixe-inclusion}
Par définition, $K \subseteq L^{\Aut(L/K)}$. L'inclusion peut être
stricte~: si $L/K$ n'est pas «~assez symétrique~», le corps fixe est
plus gros que $K$. L'égalité $K = L^{\Aut(L/K)}$ caractérise précisément
les extensions galoisiennes, comme nous le verrons.
\end{remarque}

\begin{lemme}[Lemme d'Artin]\label{lem:artin}
\index{Artin!lemme d'}
Soit $L$ un corps et $H$ un sous-groupe \emph{fini} de $\Aut(L)$. Alors
\[
  [L : L^H] = |H|.
\]
En particulier, $L/L^H$ est une extension finie.
\end{lemme}

\begin{proof}
Posons $K = L^H$ et $n = |H|$. La preuve se déroule en deux étapes.

\medskip\noindent
\textbf{Étape 1 :} $[L:K] \leq n$.

Supposons par l'absurde que $[L:K] > n$. Alors il existe $n+1$ éléments
$\alpha_1, \ldots, \alpha_{n+1} \in L$ qui sont $K$-linéairement indépendants.
Soient $\sigma_1 = \mathrm{id}, \sigma_2, \ldots, \sigma_n$ les éléments
de $H$. Considérons le système linéaire homogène
\[
  \sum_{j=1}^{n+1} x_j \, \sigma_i(\alpha_j) = 0, \quad i = 1, \ldots, n.
\]
Ce système a $n$ équations et $n+1$ inconnues à coefficients dans $L$,
donc il admet une solution non triviale $(x_1, \ldots, x_{n+1}) \in L^{n+1}
\setminus \{0\}$. Choisissons une telle solution avec un nombre minimal de
composantes non nulles~; quitte à réordonner et diviser, on peut supposer
$x_1 = 1$. Pour $i = 1$ (correspondant à $\sigma_1 = \mathrm{id}$), on
obtient $\alpha_1 + x_2 \alpha_2 + \cdots + x_{n+1}\alpha_{n+1} = 0$.
Comme les $\alpha_j$ sont $K$-linéairement indépendants, au moins un $x_j$
n'appartient pas à $K = L^H$, disons $x_2 \notin K$.

Il existe donc $\tau \in H$ avec $\tau(x_2) \neq x_2$. Appliquons $\tau$
à l'égalité $\sum_j x_j \, \sigma_i(\alpha_j) = 0$ pour tout $i$~: on
obtient $\sum_j \tau(x_j)\, (\tau\sigma_i)(\alpha_j) = 0$ pour tout $i$.
Comme $H$ est un groupe, $\{\tau\sigma_1, \ldots, \tau\sigma_n\}$ est une
permutation de $\{\sigma_1, \ldots, \sigma_n\}$, donc
$\sum_j \tau(x_j)\,\sigma_i(\alpha_j) = 0$ pour tout $i$.

En soustrayant de la relation originale~:
$\sum_{j=2}^{n+1} (x_j - \tau(x_j))\,\sigma_i(\alpha_j) = 0$ pour tout $i$
(puisque $x_1 = 1 = \tau(x_1)$). C'est une solution non triviale (car
$x_2 - \tau(x_2) \neq 0$) ayant strictement moins de composantes non
nulles, contradiction.

\medskip\noindent
\textbf{Étape 2 :} $[L:K] \geq n$.

Par la proposition~\ref{prop:majoration-aut}, $|\Aut(L/K)| \leq [L:K]$.
Or $H \leq \Aut(L/K)$ puisque tout $\sigma \in H$ fixe $K = L^H$,
donc $n = |H| \leq |\Aut(L/K)| \leq [L:K]$.

\medskip
Des deux inégalités, $[L:K] = n = |H|$.
\end{proof}

% ------------------------------------------------------------
\section{Extensions normales}\label{sec:extensions-normales}
\index{extension!normale}
% ------------------------------------------------------------

\begin{definition}[Extension normale]\label{def:extension-normale}
Une extension algébrique $L/K$ est dite \emph{normale} si tout polynôme
irréductible $f \in K[X]$ qui possède au moins une racine dans $L$ se
scinde complètement dans $L[X]$.
\end{definition}

\begin{theoreme}[Caractérisation des extensions normales finies]\label{thm:normale-splitting}
\index{corps!de décomposition}
Soit $L/K$ une extension finie. Les assertions suivantes sont équivalentes~:
\begin{enumerate}[label=(\roman*)]
  \item $L/K$ est normale.
  \item $L$ est un corps de décomposition d'un polynôme $f \in K[X]$.
  \item Pour tout plongement $\sigma \colon L \hookrightarrow \overline{K}$
    qui fixe $K$, on a $\sigma(L) = L$ (c'est-à-dire $\sigma(L) \subseteq L$).
\end{enumerate}
\end{theoreme}

\begin{proof}
\textbf{(i) $\Rightarrow$ (ii).}
Puisque $L/K$ est finie, on peut écrire $L = K(\alpha_1, \ldots, \alpha_r)$.
Pour chaque $i$, soit $f_i = \irr(\alpha_i, K)$. Par hypothèse de normalité,
chaque $f_i$ se scinde dans $L[X]$. Posons $f = f_1 \cdots f_r \in K[X]$.
Alors $f$ se scinde dans $L[X]$, et $L$ est engendré sur $K$ par un
sous-ensemble des racines de $f$ (les $\alpha_i$). Si $L'$ désigne le
corps de décomposition de $f$ sur $K$ contenu dans $L$, alors
$L' \supseteq K(\alpha_1, \ldots, \alpha_r) = L$, d'où $L' = L$.

\medskip\noindent
\textbf{(ii) $\Rightarrow$ (iii).}
Supposons que $L$ est le corps de décomposition de $f \in K[X]$. Écrivons
$f = a\prod_{i=1}^n (X - \beta_i)$ dans $L[X]$, de sorte que
$L = K(\beta_1, \ldots, \beta_n)$. Soit $\sigma \colon L \hookrightarrow
\overline{K}$ un $K$-plongement. Pour chaque $i$, $\sigma(\beta_i)$ est
une racine de $f$ dans $\overline{K}$. Or toutes les racines de $f$ sont
déjà dans $L$ (car $f$ se scinde dans $L$), donc $\sigma(\beta_i) \in L$.
Puisque $L = K(\beta_1, \ldots, \beta_n)$, on en déduit
$\sigma(L) \subseteq L$. Comme $\sigma$ est injectif et $L$ est de
dimension finie sur $K$, c'est un automorphisme de $L$, d'où
$\sigma(L) = L$.

\medskip\noindent
\textbf{(iii) $\Rightarrow$ (i).}
Soit $g \in K[X]$ irréductible ayant une racine $\alpha \in L$. Soit
$\beta$ une autre racine de $g$ dans $\overline{K}$. Puisque $g$ est
irréductible, il existe un $K$-isomorphisme $\tau \colon K(\alpha) \to
K(\beta)$ envoyant $\alpha$ sur $\beta$. Ce $\tau$ se prolonge en un
plongement $\sigma \colon L \hookrightarrow \overline{K}$ (par le théorème
de prolongement). Par hypothèse, $\sigma(L) = L$, donc
$\beta = \sigma(\alpha) \in L$. Ainsi toutes les racines de $g$ dans
$\overline{K}$ sont dans $L$, et $g$ se scinde dans $L[X]$.
\end{proof}

% ------------------------------------------------------------
\section{Extensions galoisiennes}\label{sec:def-galoisienne}
\index{extension!galoisienne}
% ------------------------------------------------------------

\begin{definition}[Extension galoisienne]\label{def:extension-galoisienne}
Une extension algébrique $L/K$ est dite \emph{galoisienne} si elle est
à la fois \emph{normale}\index{extension!normale} et
\emph{séparable}\index{extension!séparable}.
\end{definition}

\begin{definition}[Groupe de Galois]\label{def:groupe-galois}
\index{groupe!de Galois}
Si $L/K$ est une extension galoisienne, le \emph{groupe de Galois} de
$L/K$ est
\[
  \Gal(L/K) = \Aut(L/K).
\]
\end{definition}

\begin{theoreme}[Caractérisations des extensions galoisiennes finies]\label{thm:caract-galoisienne}
Soit $L/K$ une extension finie. Les assertions suivantes sont équivalentes~:
\begin{enumerate}[label=(\roman*)]
  \item $L/K$ est galoisienne (normale et séparable).
  \item $L$ est le corps de décomposition d'un polynôme séparable
    $f \in K[X]$.
  \item $|\Aut(L/K)| = [L:K]$.
  \item $K = L^{\Aut(L/K)}$, c'est-à-dire $K$ est le corps fixe du
    groupe d'automorphismes.
\end{enumerate}
\end{theoreme}

\begin{proof}
\textbf{(i) $\Rightarrow$ (ii).}
Écrivons $L = K(\alpha_1, \ldots, \alpha_r)$. Pour chaque $i$, le polynôme
$f_i = \irr(\alpha_i, K)$ est séparable (car $L/K$ est séparable) et se
scinde dans $L$ (car $L/K$ est normale). Posons $f = \mathrm{ppcm}(f_1,
\ldots, f_r)$~: c'est un polynôme séparable de $K[X]$ qui se scinde dans $L$.
Les racines de $f$ engendrent $L$ sur $K$, donc $L$ est le corps de
décomposition de $f$.

\medskip\noindent
\textbf{(ii) $\Rightarrow$ (iii).}
Supposons que $L$ est le corps de décomposition d'un polynôme séparable
$f \in K[X]$. On procède par récurrence sur $[L:K]$. Si $[L:K] = 1$, le
résultat est trivial. Sinon, soit $g$ un facteur irréductible de $f$ de
degré $d \geq 2$, et soit $\alpha$ une racine de $g$ dans $L$. Le polynôme
$g$ est séparable, donc il a $d$ racines distinctes $\alpha = \alpha_1,
\alpha_2, \ldots, \alpha_d$ dans $L$. Pour chaque $j$, il existe un
$K$-isomorphisme $\tau_j \colon K(\alpha) \to K(\alpha_j)$ envoyant $\alpha$
sur $\alpha_j$. Or $L$ est le corps de décomposition de $f$ sur $K(\alpha)$
comme sur $K(\alpha_j)$~; par l'unicité du corps de décomposition, $\tau_j$
se prolonge en un automorphisme $\sigma_j$ de $L$ fixant $K$. Ces $\sigma_j$
sont deux à deux distincts sur $K(\alpha)$.

Par hypothèse de récurrence appliquée à $L/K(\alpha)$ (qui est le corps de
décomposition de $f$ sur $K(\alpha)$, $f$ restant séparable), on a
$|\Aut(L/K(\alpha))| = [L:K(\alpha)]$. Les classes
$\sigma_j \cdot \Aut(L/K(\alpha))$ pour $j = 1, \ldots, d$ sont deux à deux
distinctes (car les $\sigma_j$ diffèrent sur $\alpha$), donc
\[
  |\Aut(L/K)| \geq d \cdot |\Aut(L/K(\alpha))| = d \cdot [L:K(\alpha)]
  = [K(\alpha):K] \cdot [L:K(\alpha)] = [L:K].
\]
Combiné avec l'inégalité $|\Aut(L/K)| \leq [L:K]$ (proposition~\ref{prop:majoration-aut}), on obtient l'égalité.

\medskip\noindent
\textbf{(iii) $\Rightarrow$ (iv).}
Posons $H = \Aut(L/K)$ et $E = L^H$. On a $K \subseteq E$, donc
$[L:K] = [L:E][E:K]$. Par le lemme d'Artin
(lemme~\ref{lem:artin}), $[L:E] = |H| = |\Aut(L/K)| = [L:K]$.
Donc $[E:K] = 1$, c'est-à-dire $E = K$.

\medskip\noindent
\textbf{(iv) $\Rightarrow$ (i).}
Posons $H = \Aut(L/K)$ et supposons $K = L^H$. Par le lemme d'Artin,
$[L:K] = |H|$. Il faut montrer que $L/K$ est normale et séparable.

\emph{Séparabilité.} Soit $\alpha \in L$ et $\mu = \irr(\alpha, K)$.
Les éléments $\sigma(\alpha)$ pour $\sigma \in H$ sont des racines de $\mu$
(car $\mu \in K[X]$ et $\sigma$ fixe $K$). Les éléments distincts parmi
les $\sigma(\alpha)$ fournissent des racines distinctes de $\mu$, donc
$\mu$ est séparable.

\emph{Normalité.} Avec les notations précédentes, le polynôme
$g = \prod_{\beta} (X - \beta)$ où $\beta$ parcourt les éléments
\emph{distincts} de $\{\sigma(\alpha) \mid \sigma \in H\}$ est invariant
par tout $\sigma \in H$ (car $H$ permute les $\sigma(\alpha)$). Ses
coefficients sont donc dans $L^H = K$, d'où $g \in K[X]$. Comme $g$
divise $\mu$ (toutes ses racines sont des racines de $\mu$) et $\mu$
divise $g$ (car $\mu$ est irréductible et $g(\alpha) = 0$), on a
$g = \mu$. Ainsi $\mu$ se scinde dans $L$.
\end{proof}

% ------------------------------------------------------------
\section{Exemples fondamentaux}\label{sec:exemples-galois}
% ------------------------------------------------------------

\begin{exemple}[$\Gal(\Q(\sqrt{2})/\Q) \cong \Z/2\Z$]\label{ex:gal-sqrt2}
\index{groupe!de Galois!de $\Q(\sqrt{2})/\Q$}
L'extension $\Q(\sqrt{2})/\Q$ est le corps de décomposition de $X^2 - 2$,
qui est séparable (car $\car(\Q) = 0$). C'est donc une extension
galoisienne de degré $2$. Tout $\sigma \in \Gal(\Q(\sqrt{2})/\Q)$ vérifie
$\sigma(\sqrt{2})^2 = \sigma(2) = 2$, donc $\sigma(\sqrt{2}) = \pm\sqrt{2}$.
Les deux choix définissent deux automorphismes distincts~:
\[
  \mathrm{id} \colon \sqrt{2} \mapsto \sqrt{2}, \qquad
  \sigma \colon \sqrt{2} \mapsto -\sqrt{2}.
\]
On a $\sigma^2 = \mathrm{id}$, donc $\Gal(\Q(\sqrt{2})/\Q)
= \{\mathrm{id}, \sigma\} \cong \Z/2\Z$.
\end{exemple}

\begin{exemple}[{$\Gal(\Q(\sqrt[3]{2}, \omega)/\Q) \cong S_3$}]\label{ex:gal-cube-root-2}
\index{groupe!de Galois!de $\Q(\sqrt[3]{2},\omega)/\Q$}
Soit $\omega = e^{2\pi i/3} = \frac{-1+i\sqrt{3}}{2}$, racine primitive
cubique de l'unité. Le polynôme $X^3 - 2$ est irréductible sur $\Q$ (par
Eisenstein avec $p = 2$) et ses racines sont
$\alpha = \sqrt[3]{2}$, $\omega\alpha$, $\omega^2\alpha$.
Le corps de décomposition est $L = \Q(\alpha, \omega)$.

\emph{Calcul du degré.} On a $[\Q(\alpha):\Q] = 3$. Le polynôme minimal
de $\omega$ sur $\Q(\alpha)$ est $X^2 + X + 1$ (qui est irréductible sur
$\Q$ et reste irréductible sur $\Q(\alpha)$ puisque $\Q(\alpha) \subset \R$
et $\omega \notin \R$). Donc $[L:\Q(\alpha)] = 2$ et
$[L:\Q] = 3 \cdot 2 = 6$.

L'extension $L/\Q$ est galoisienne (corps de décomposition d'un polynôme
séparable), donc $|\Gal(L/\Q)| = 6$.

\emph{Détermination des automorphismes.} Tout $\sigma \in \Gal(L/\Q)$ est
déterminé par $\sigma(\alpha)$ et $\sigma(\omega)$~:
\begin{itemize}
  \item $\sigma(\alpha) \in \{\alpha, \omega\alpha, \omega^2\alpha\}$
    (racines de $X^3-2$),
  \item $\sigma(\omega) \in \{\omega, \omega^2\}$
    (racines de $X^2+X+1$).
\end{itemize}
Cela donne $3 \times 2 = 6$ possibilités, et chacune se réalise~:

\begin{center}
\renewcommand{\arraystretch}{1.2}
\begin{tabular}{c|cc}
  $\sigma$ & $\sigma(\alpha)$ & $\sigma(\omega)$ \\
  \hline
  $\mathrm{id}$ & $\alpha$ & $\omega$ \\
  $\rho$ & $\omega\alpha$ & $\omega$ \\
  $\rho^2$ & $\omega^2\alpha$ & $\omega$ \\
  $\tau$ & $\alpha$ & $\omega^2$ \\
  $\rho\tau$ & $\omega\alpha$ & $\omega^2$ \\
  $\rho^2\tau$ & $\omega^2\alpha$ & $\omega^2$
\end{tabular}
\end{center}
On vérifie que $\rho$ est d'ordre $3$ (permutation cyclique des racines),
$\tau$ est d'ordre $2$ (conjugaison complexe), et $\tau\rho\tau = \rho^{-1}$.
Le groupe $G = \langle \rho, \tau \mid \rho^3 = \tau^2 = 1,\
\tau\rho\tau = \rho^{-1}\rangle$ est isomorphe à $S_3$\index{$S_3$}.

Si l'on identifie les racines $\alpha, \omega\alpha, \omega^2\alpha$ aux
éléments $\{1, 2, 3\}$, l'action de $G$ par permutation des racines fournit
un isomorphisme $\Gal(L/\Q) \cong S_3$.
\end{exemple}

\begin{exemple}[$\Gal(\F_{p^n}/\F_p) \cong \Z/n\Z$]\label{ex:gal-corps-finis}
\index{Frobenius!automorphisme de}
\index{corps!fini!groupe de Galois}
Soit $p$ un nombre premier et $n \geq 1$. Le corps $\F_{p^n}$ est le corps
de décomposition de $X^{p^n} - X$ sur $\F_p$ ; ce polynôme est séparable
(sa dérivée est $-1 \neq 0$). Donc $\F_{p^n}/\F_p$ est galoisienne de
degré $n$, et $|\Gal(\F_{p^n}/\F_p)| = n$.

Considérons l'automorphisme de \emph{Frobenius}
\[
  \varphi \colon \F_{p^n} \to \F_{p^n}, \quad x \mapsto x^p.
\]
C'est bien un automorphisme de corps (l'identité $(a+b)^p = a^p + b^p$
en caractéristique $p$ assure que c'est un homomorphisme d'anneaux, et
l'injectivité entraîne la surjectivité en dimension finie).
De plus, $\varphi \in \Gal(\F_{p^n}/\F_p)$ car $a^p = a$ pour tout
$a \in \F_p$.

\emph{Affirmation :} $\varphi$ est d'ordre exactement $n$.

En effet, $\varphi^k(x) = x^{p^k}$ pour tout $x \in \F_{p^n}$.
Si $\varphi^k = \mathrm{id}$, alors $x^{p^k} = x$ pour tout
$x \in \F_{p^n}$, ce qui signifie que tout élément de $\F_{p^n}$ est
racine de $X^{p^k} - X$. Ce polynôme a au plus $p^k$ racines, donc
$p^n \leq p^k$, d'où $k \geq n$. Inversement, $\varphi^n(x) = x^{p^n} = x$
pour tout $x \in \F_{p^n}$ (car $\F_{p^n}^* = \{x : x^{p^n-1} = 1\}$).
Donc l'ordre de $\varphi$ est $n$.

Ainsi $\langle\varphi\rangle$ est un sous-groupe cyclique d'ordre $n$ de
$\Gal(\F_{p^n}/\F_p)$, qui est lui-même d'ordre $n$. On en déduit
\[
  \Gal(\F_{p^n}/\F_p) = \langle\varphi\rangle \cong \Z/n\Z.
\]
\end{exemple}

\begin{exemple}[Extensions cyclotomiques]\label{ex:gal-cyclotomique}
\index{extension!cyclotomique}
\index{polynôme!cyclotomique}
Soit $n \geq 1$ un entier et $\zeta_n = e^{2\pi i/n}$ une racine primitive
$n$-ième de l'unité. Le $n$-ième polynôme cyclotomique est
\[
  \Phi_n(X) = \prod_{\substack{1 \leq k \leq n \\ \pgcd(k,n)=1}}
  (X - \zeta_n^k) \in \Z[X].
\]
C'est un polynôme irréductible sur $\Q$ (résultat classique). Son degré
est $\varphi(n)$, la fonction indicatrice d'Euler\index{Euler!indicatrice d'}.

L'extension $\Q(\zeta_n)/\Q$ est le corps de décomposition de $X^n - 1$,
qui est séparable sur $\Q$. C'est donc une extension galoisienne de degré
$\varphi(n)$.

Tout $\sigma \in \Gal(\Q(\zeta_n)/\Q)$ est déterminé par $\sigma(\zeta_n)$,
qui doit être une racine primitive $n$-ième de l'unité~: $\sigma(\zeta_n)
= \zeta_n^{a_\sigma}$ pour un unique $a_\sigma \in (\Z/n\Z)^*$. L'application
\[
  \chi \colon \Gal(\Q(\zeta_n)/\Q) \to (\Z/n\Z)^*, \quad
  \sigma \mapsto a_\sigma
\]
est un homomorphisme de groupes injectif (car $\sigma$ est déterminé par
$a_\sigma$). Puisque les deux groupes ont même cardinal $\varphi(n)$,
c'est un isomorphisme~:
\[
  \Gal(\Q(\zeta_n)/\Q) \cong (\Z/n\Z)^*.
\]
\end{exemple}

% ------------------------------------------------------------
\section{Tour d'extensions et diagrammes}\label{sec:tours-extensions}
% ------------------------------------------------------------

Le diagramme suivant illustre une tour d'extensions avec les groupes de Galois
associés.

\begin{center}
\begin{tikzcd}[row sep=1.5cm, column sep=0.5cm]
  L \arrow[d, dash, "{|H|}"'] \arrow[dd, dash, bend left=50,
    "{|G| = [L:K]}" {right, pos=0.5}] \\
  L^H \arrow[d, dash, "{[G:H]}"'] \\
  K
\end{tikzcd}
\end{center}

\noindent
Ici, $G = \Gal(L/K)$ et $H \leq G$. Le lemme d'Artin donne $[L:L^H] = |H|$,
et la formule des degrés entraîne $[L^H:K] = [L:K]/[L:L^H] = |G|/|H| = [G:H]$.

\begin{center}
\begin{tikzcd}[row sep=1.2cm]
  \Q(\sqrt[3]{2}, \omega) \arrow[d, dash, "2"'] \arrow[dr, dash, "3"]
    \arrow[dd, dash, bend right=60, "6"'] \\
  \Q(\sqrt[3]{2}) \arrow[d, dash, "3"'] & \Q(\omega) \arrow[dl, dash, "2"] \\
  \Q
\end{tikzcd}
\end{center}

% ------------------------------------------------------------
\section{Exercices}\label{sec:exercices-chap11}
% ------------------------------------------------------------

\begin{exercice}[Automorphismes de $\Q(\sqrt{2},\sqrt{3})$]\label{exo:aut-sqrt2-sqrt3}
Montrer que $\Q(\sqrt{2},\sqrt{3})/\Q$ est galoisienne et que
$\Gal(\Q(\sqrt{2},\sqrt{3})/\Q) \cong (\Z/2\Z)^2$.
\emph{Indication :} montrer que $[\Q(\sqrt{2},\sqrt{3}):\Q] = 4$ et
exhiber quatre automorphismes.
\end{exercice}

\begin{exercice}[Corps fixe explicite]\label{exo:corps-fixe-explicite}
Soit $L = \Q(\sqrt[4]{2}, i)$, qui est le corps de décomposition de $X^4 - 2$
sur $\Q$. Montrer que $[L:\Q] = 8$ et que $\Gal(L/\Q) \cong D_4$ (le groupe
diédral d'ordre $8$). Déterminer le corps fixe du sous-groupe
$\langle\sigma\rangle$ où $\sigma \colon \sqrt[4]{2} \mapsto i\sqrt[4]{2}$,
$i \mapsto i$.
\end{exercice}

\begin{exercice}[Frobenius et sous-corps]\label{exo:frobenius-sous-corps}
Soit $\varphi$ le Frobenius de $\F_{p^{12}}/\F_p$. Déterminer tous les
sous-groupes de $\Gal(\F_{p^{12}}/\F_p) = \langle\varphi\rangle \cong
\Z/12\Z$ et les sous-corps correspondants.
\end{exercice}

\begin{exercice}[Extension non galoisienne]\label{exo:ext-non-galoisienne}
Montrer que $\Q(\sqrt[3]{2})/\Q$ n'est pas une extension galoisienne.
Calculer $|\Aut(\Q(\sqrt[3]{2})/\Q)|$.
\end{exercice}

\begin{exercice}[Groupe de Galois d'un polynôme cyclotomique]\label{exo:gal-cyclotomique-12}
Déterminer $\Gal(\Q(\zeta_{12})/\Q)$ et identifier ce groupe.
\end{exercice}

\begin{exercice}[Produit de groupes de Galois]\label{exo:produit-galois}
Soient $L_1/K$ et $L_2/K$ deux extensions galoisiennes finies contenues
dans une même extension $\Omega/K$, avec $L_1 \cap L_2 = K$. Montrer que
$L_1 L_2 / K$ est galoisienne et que
$\Gal(L_1 L_2/K) \cong \Gal(L_1/K) \times \Gal(L_2/K)$.
\end{exercice}

\begin{exercice}[Corps fixe et trace]\label{exo:corps-fixe-trace}
Soit $L/K$ une extension galoisienne finie de groupe $G = \Gal(L/K)$.
La \emph{trace} est l'application $K$-linéaire
$\mathrm{Tr}_{L/K} \colon L \to K$, $\alpha \mapsto \sum_{\sigma \in G}
\sigma(\alpha)$.
\begin{enumerate}[label=(\alph*)]
  \item Vérifier que $\mathrm{Tr}_{L/K}(\alpha) \in K$ pour tout
    $\alpha \in L$.
  \item Montrer que $\mathrm{Tr}_{L/K}$ est surjective.
    \emph{Indication :} utiliser l'indépendance des caractères (lemme de
    Dedekind).
\end{enumerate}
\end{exercice}

\begin{exercice}[Lemme de Dedekind]\label{exo:lemme-dedekind}
\index{Dedekind!lemme de}
Soient $L, L'$ deux corps et $\sigma_1, \ldots, \sigma_n \colon L \to L'$
des homomorphismes de corps deux à deux distincts. Montrer que
$\sigma_1, \ldots, \sigma_n$ sont $L'$-linéairement indépendants dans
l'espace des applications $L \to L'$.
\end{exercice}

\begin{exercice}[Clôture galoisienne]\label{exo:cloture-galoisienne}
Soit $L/K$ une extension finie séparable. Montrer qu'il existe une plus petite
extension galoisienne $N/K$ contenant $L$ (dans une clôture algébrique fixée).
Calculer $N$ pour $L = \Q(\sqrt[3]{2})$.
\end{exercice}

\begin{exercice}[Automorphismes en caractéristique $p$]\label{exo:aut-car-p}
Soit $K = \F_p(t)$ le corps des fractions rationnelles sur $\F_p$ et
$L = K(\alpha)$ où $\alpha^p - \alpha = t$. Montrer que $L/K$ est une
extension galoisienne de groupe $\Z/p\Z$.
\emph{Indication :} les racines du polynôme $X^p - X - t$ sont
$\alpha, \alpha+1, \ldots, \alpha+(p-1)$.
\end{exercice}

\bigskip
\begin{center}
\fbox{\parbox{0.85\textwidth}{
\textbf{Résumé du chapitre~\ref{chap:extensions-galoisiennes}.}
Une extension $L/K$ est galoisienne si et seulement si elle est normale et
séparable, si et seulement si $|\Aut(L/K)| = [L:K]$, si et seulement si
$K = L^{\Aut(L/K)}$. Le groupe $\Gal(L/K) = \Aut(L/K)$ encode les symétries
de l'extension. Les exemples clefs sont~: $\Gal(\Q(\sqrt{d})/\Q) \cong
\Z/2\Z$, $\Gal(\Q(\sqrt[3]{2},\omega)/\Q) \cong S_3$,
$\Gal(\F_{p^n}/\F_p) \cong \Z/n\Z$,
$\Gal(\Q(\zeta_n)/\Q) \cong (\Z/n\Z)^*$.
Le lemme d'Artin ($[L:L^H] = |H|$) est un outil central.
}}
\end{center}

% ============================================================
%  CHAPITRE 12 : THÉORIE DE GALOIS — THÉORÈME FONDAMENTAL
% ============================================================
\chapter{Théorie de Galois --- Théorème fondamental}\label{chap:theoreme-fondamental}
\index{théorème!fondamental de la théorie de Galois}

\bigskip
\noindent
Le théorème fondamental de la théorie de Galois établit une correspondance
bijective entre les sous-groupes du groupe de Galois et les corps
intermédiaires d'une extension galoisienne finie. Cette correspondance, qui
\emph{renverse l'ordre d'inclusion}, constitue l'un des plus beaux résultats
de l'algèbre~: elle traduit des questions sur les corps (objets a priori
compliqués) en questions sur les groupes finis (objets plus maniables et mieux
compris). C'est un paradigme que l'on retrouve dans de nombreux domaines des
mathématiques --- de la topologie (revêtements et groupe fondamental) à la
géométrie algébrique (correspondance de Galois pour les revêtements étales).

% ------------------------------------------------------------
\section{Énoncé et preuve du théorème fondamental}\label{sec:enonce-preuve-FTGT}
% ------------------------------------------------------------

\begin{theoreme}[Théorème fondamental de la théorie de Galois]\label{thm:fondamental-galois}
\index{correspondance de Galois}
Soit $L/K$ une extension galoisienne finie de groupe de Galois
$G = \Gal(L/K)$. Notons
\[
  \mathcal{E} = \{E \mid K \subseteq E \subseteq L,\ E \text{ sous-corps}\},
  \qquad
  \mathcal{H} = \{H \mid H \leq G\}.
\]
Les applications
\[
  \Phi \colon \mathcal{E} \to \mathcal{H}, \quad E \mapsto \Gal(L/E),
  \qquad
  \Psi \colon \mathcal{H} \to \mathcal{E}, \quad H \mapsto L^H
\]
sont des bijections réciproques l'une de l'autre, renversant l'inclusion.
De plus~:
\begin{enumerate}[label=(\roman*)]
  \item \textbf{Bijection.} $\Psi \circ \Phi = \mathrm{id}_{\mathcal{E}}$
    et $\Phi \circ \Psi = \mathrm{id}_{\mathcal{H}}$.
  \item \textbf{Degrés et indices.} Pour tout $E \in \mathcal{E}$ et
    $H = \Gal(L/E)$~:
    \[
      [L:E] = |H|, \qquad [E:K] = [G:H].
    \]
  \item \textbf{Normalité.} $E/K$ est une extension normale (et donc
    galoisienne, puisque séparable) si et seulement si $\Gal(L/E)$ est
    un sous-groupe distingué de $G$. Dans ce cas,
    \[
      \Gal(E/K) \cong G / \Gal(L/E).
    \]
\end{enumerate}
\end{theoreme}

\begin{proof}
Nous découpons la preuve en trois parties correspondant aux trois assertions.

\medskip\noindent
\textbf{Partie (i) : la bijection.}

\emph{$\Psi \circ \Phi = \mathrm{id}_{\mathcal{E}}$.}
Soit $E \in \mathcal{E}$. On doit montrer $L^{\Gal(L/E)} = E$.
L'extension $L/E$ est galoisienne (elle est normale car $L$ est le corps de
décomposition d'un polynôme séparable $f \in K[X] \subset E[X]$, et
séparable car sous-extension d'une extension séparable). Par le
théorème~\ref{thm:caract-galoisienne}~(iv) appliqué à $L/E$, on a
$E = L^{\Aut(L/E)} = L^{\Gal(L/E)}$.

\emph{$\Phi \circ \Psi = \mathrm{id}_{\mathcal{H}}$.}
Soit $H \leq G$. On doit montrer $\Gal(L/L^H) = H$. Par le lemme d'Artin
(lemme~\ref{lem:artin}), $[L:L^H] = |H|$. De plus, $L/L^H$ est galoisienne
(même argument que ci-dessus), donc $|\Gal(L/L^H)| = [L:L^H] = |H|$.
Or $H \leq \Gal(L/L^H)$ (tout élément de $H$ fixe $L^H$ par définition).
Un sous-groupe d'un groupe fini ayant le même cardinal que le groupe est le
groupe entier, donc $\Gal(L/L^H) = H$.

\medskip\noindent
\textbf{Partie (ii) : les degrés.}

Pour $E \in \mathcal{E}$ et $H = \Gal(L/E)$, l'extension $L/E$ est
galoisienne, donc $[L:E] = |\Gal(L/E)| = |H|$. La formule des degrés donne
\[
  [E:K] = \frac{[L:K]}{[L:E]} = \frac{|G|}{|H|} = [G:H].
\]

\medskip\noindent
\textbf{Partie (iii) : le critère de normalité.}

($\Rightarrow$) Supposons $E/K$ normale. Soit $\sigma \in G$ et
$\tau \in H = \Gal(L/E)$. On doit montrer $\sigma\tau\sigma^{-1} \in H$,
c'est-à-dire que $\sigma\tau\sigma^{-1}$ fixe $E$. Soit $\alpha \in E$.
Puisque $\sigma^{-1}$ est un $K$-automorphisme de $L$ et $E/K$ est
normale, $\sigma^{-1}$ envoie $E$ dans $E$ (théorème~\ref{thm:normale-splitting}~(iii)
appliqué à tout $K$-plongement). Donc $\sigma^{-1}(\alpha) \in E$, puis
$\tau(\sigma^{-1}(\alpha)) = \sigma^{-1}(\alpha)$ (car $\tau$ fixe $E$),
et enfin $\sigma\tau\sigma^{-1}(\alpha) = \alpha$. Ainsi $H \trianglelefteq G$.

Définissons $\rho \colon G \to \Gal(E/K)$ par $\rho(\sigma) = \sigma|_E$.
C'est bien défini car $\sigma(E) = E$ ($E/K$ normale). C'est clairement un
homomorphisme de groupes. Son noyau est
$\Ker(\rho) = \{\sigma \in G \mid \sigma|_E = \mathrm{id}_E\} = \Gal(L/E) = H$.
Il reste à montrer la surjectivité. Soit $\phi \in \Gal(E/K)$~; puisque
$L$ est le corps de décomposition d'un polynôme $f \in K[X]$ et
$E \subseteq L$, l'automorphisme $\phi$ de $E$ se prolonge en un automorphisme
$\sigma$ de $L$ (par le théorème de prolongement des isomorphismes aux corps
de décomposition). Alors $\sigma \in G$ et $\rho(\sigma) = \phi$.

Par le premier théorème d'isomorphisme,
$\Gal(E/K) \cong G/H = G/\Gal(L/E)$.

($\Leftarrow$) Supposons $H = \Gal(L/E) \trianglelefteq G$. Soit
$\alpha \in E$ et $\mu = \irr(\alpha, K)$. Soit $\beta$ une racine de $\mu$
dans $L$ (qui contient toutes les racines de $\mu$ puisque $L/K$ est
normale). Il existe $\sigma \in G$ avec $\sigma(\alpha) = \beta$. Pour tout
$\tau \in H$, on a $\sigma^{-1}\tau\sigma \in H$ (car $H$ est distingué),
donc $\sigma^{-1}\tau\sigma(\alpha) = \alpha$, c'est-à-dire
$\tau(\sigma(\alpha)) = \sigma(\alpha)$, soit $\tau(\beta) = \beta$.
Comme ceci vaut pour tout $\tau \in H = \Gal(L/E)$, on en déduit
$\beta \in L^H = E$. Donc toutes les racines de $\mu$ dans $L$ sont dans $E$,
et $\mu$ se scinde dans $E[X]$. Ainsi $E/K$ est normale.
\end{proof}

% ------------------------------------------------------------
\section{Exemples de la correspondance de Galois}\label{sec:exemples-correspondance}
% ------------------------------------------------------------

\subsection{$\Q(\sqrt{2},\sqrt{3})/\Q$ : groupe de Klein $V_4$}\label{ssec:V4}
\index{groupe!de Klein $V_4$}

Posons $L = \Q(\sqrt{2},\sqrt{3})$. Le polynôme $(X^2-2)(X^2-3)$ est
séparable et se scinde dans $L$, donc $L/\Q$ est galoisienne. On a
$[L:\Q] = 4$ et $G = \Gal(L/\Q) = \{\mathrm{id}, \sigma, \tau, \sigma\tau\}$
où~:
\[
  \sigma \colon \sqrt{2} \mapsto -\sqrt{2},\ \sqrt{3} \mapsto \sqrt{3},
  \qquad
  \tau \colon \sqrt{2} \mapsto \sqrt{2},\ \sqrt{3} \mapsto -\sqrt{3}.
\]
On a $\sigma^2 = \tau^2 = (\sigma\tau)^2 = \mathrm{id}$, donc
$G \cong V_4 = (\Z/2\Z)^2$.

Les sous-groupes de $G$ sont~:
\[
  \{e\},\quad \langle\sigma\rangle,\quad \langle\tau\rangle,\quad
  \langle\sigma\tau\rangle,\quad G.
\]
Les corps intermédiaires correspondants (par $H \mapsto L^H$) sont~:
\begin{align*}
  L^{\{e\}} &= L = \Q(\sqrt{2},\sqrt{3}), \\
  L^{\langle\sigma\rangle} &= \Q(\sqrt{3}) &
    &(\text{éléments fixés par }\sigma), \\
  L^{\langle\tau\rangle} &= \Q(\sqrt{2}) &
    &(\text{éléments fixés par }\tau), \\
  L^{\langle\sigma\tau\rangle} &= \Q(\sqrt{6}) &
    &(\text{car }\sigma\tau(\sqrt{6}) = (-\sqrt{2})(-\sqrt{3}) = \sqrt{6}), \\
  L^G &= \Q.
\end{align*}

Tous les sous-groupes de $V_4$ sont distingués (le groupe est abélien), donc
toutes les sous-extensions sont normales (et galoisiennes) sur $\Q$.

\begin{center}
\begin{tikzcd}[row sep=1.2cm, column sep=0.8cm]
  & L = \Q(\sqrt{2},\sqrt{3}) \arrow[dl, dash, "2"']
    \arrow[d, dash, "2"] \arrow[dr, dash, "2"] & \\
  \Q(\sqrt{3}) \arrow[dr, dash, "2"'] &
  \Q(\sqrt{6}) \arrow[d, dash, "2"] &
  \Q(\sqrt{2}) \arrow[dl, dash, "2"] \\
  & \Q &
\end{tikzcd}
$\qquad\longleftrightarrow\qquad$
\begin{tikzcd}[row sep=1.2cm, column sep=0.6cm]
  & \{e\} \arrow[dl, dash] \arrow[d, dash] \arrow[dr, dash] & \\
  \langle\sigma\rangle \arrow[dr, dash] &
  \langle\sigma\tau\rangle \arrow[d, dash] &
  \langle\tau\rangle \arrow[dl, dash] \\
  & G = V_4 &
\end{tikzcd}
\end{center}

\noindent
\emph{On observe bien le renversement de l'inclusion}~: le plus gros corps
correspond au plus petit groupe, et réciproquement.

\subsection{$\Q(\sqrt[3]{2},\omega)/\Q$ : groupe $S_3$}\label{ssec:S3}

Reprenons l'extension $L = \Q(\sqrt[3]{2},\omega)/\Q$ de
l'exemple~\ref{ex:gal-cube-root-2}, avec $G = \Gal(L/\Q) \cong S_3$.
Les sous-groupes de $S_3$ sont~:
\begin{itemize}
  \item $\{e\}$ (trivial),
  \item $\langle\rho\rangle \cong \Z/3\Z = A_3$ (unique sous-groupe d'indice $2$,
    distingué),
  \item $\langle\tau\rangle$, $\langle\rho\tau\rangle$,
    $\langle\rho^2\tau\rangle$ (trois sous-groupes d'ordre $2$, conjugués,
    non distingués),
  \item $G = S_3$.
\end{itemize}

Les corps intermédiaires correspondants~:
\begin{align*}
  L^{\{e\}} &= L, \\
  L^{\langle\rho\rangle} &= \Q(\omega)
    &(\text{car }\rho\text{ fixe }\omega\text{ et permute les }\sqrt[3]{2}\omega^k), \\
  L^{\langle\tau\rangle} &= \Q(\sqrt[3]{2})
    &(\text{car }\tau\text{ fixe }\sqrt[3]{2}\text{ et envoie }\omega
    \text{ sur }\omega^2), \\
  L^{\langle\rho\tau\rangle} &= \Q(\omega\sqrt[3]{2}), \\
  L^{\langle\rho^2\tau\rangle} &= \Q(\omega^2\sqrt[3]{2}), \\
  L^G &= \Q.
\end{align*}

Le seul sous-groupe distingué non trivial propre de $S_3$ est
$A_3 = \langle\rho\rangle$, donc le seul sous-corps intermédiaire
$E$ tel que $E/\Q$ soit normale est $\Q(\omega)$. On a bien
$\Gal(\Q(\omega)/\Q) \cong S_3/A_3 \cong \Z/2\Z$.

Les extensions $\Q(\sqrt[3]{2})/\Q$, $\Q(\omega\sqrt[3]{2})/\Q$ et
$\Q(\omega^2\sqrt[3]{2})/\Q$ ne sont \emph{pas} normales sur $\Q$ (les
sous-groupes correspondants ne sont pas distingués).

\begin{center}
\begin{tikzcd}[row sep=1.3cm, column sep=0.4cm]
  & & L \arrow[dll, dash, "2"'] \arrow[dl, dash, "3"']
    \arrow[d, dash, "3"] \arrow[dr, dash, "3"] & & \\
  \Q(\omega) \arrow[drr, dash, "3"'] &
  \Q(\sqrt[3]{2}) \arrow[dr, dash, "2"'] &
  \Q(\omega\sqrt[3]{2}) \arrow[d, dash, "2"] &
  \Q(\omega^2\sqrt[3]{2}) \arrow[dl, dash, "2"] & \\
  & & \Q & &
\end{tikzcd}
$\qquad\longleftrightarrow\qquad$
\begin{tikzcd}[row sep=1.3cm, column sep=0.1cm]
  & & \{e\} \arrow[dll, dash] \arrow[dl, dash]
    \arrow[d, dash] \arrow[dr, dash] & & \\
  \langle\rho\rangle \arrow[drr, dash] &
  \langle\tau\rangle \arrow[dr, dash] &
  \langle\rho\tau\rangle \arrow[d, dash] &
  \langle\rho^2\tau\rangle \arrow[dl, dash] & \\
  & & S_3 & &
\end{tikzcd}
\end{center}

\subsection{$\F_{p^{12}}/\F_p$ : groupe cyclique $\Z/12\Z$}\label{ssec:Fp12}

Le groupe $\Gal(\F_{p^{12}}/\F_p) = \langle\varphi\rangle \cong \Z/12\Z$
est cyclique. Ses sous-groupes sont les $\langle\varphi^d\rangle$ pour
$d \mid 12$, d'ordre $12/d$. Les diviseurs de $12$ sont
$1, 2, 3, 4, 6, 12$. Le sous-corps correspondant à
$\langle\varphi^d\rangle$ est $\F_{p^d}$ (car c'est le corps fixe
de l'automorphisme $x \mapsto x^{p^d}$, qui a exactement $p^d$ points
fixes). On obtient le treillis~:

\begin{center}
\begin{tikzcd}[row sep=1cm, column sep=0.3cm]
  & & \F_{p^{12}} \arrow[dl, dash] \arrow[dr, dash] & & \\
  & \F_{p^6} \arrow[dl, dash] \arrow[d, dash] & &
    \F_{p^4} \arrow[dll, dash] & \\
  \F_{p^3} \arrow[dr, dash] & \F_{p^2} \arrow[d, dash] &
    & & \\
  & \F_p & & &
\end{tikzcd}
\end{center}

\noindent
Tous les sous-groupes d'un groupe cyclique sont distingués, donc tous les
sous-corps sont des extensions normales (et galoisiennes) de $\F_p$.

\begin{remarque}[Correspondance et divisibilité]\label{rem:treillis-corps-finis}
Pour les corps finis, le treillis des sous-corps de $\F_{p^n}$ contenant
$\F_p$ est isomorphe au treillis des diviseurs de $n$ (ordonné par la
divisibilité), lui-même anti-isomorphe au treillis des sous-groupes de
$\Z/n\Z$.
\end{remarque}

% ------------------------------------------------------------
\section{Théorème de l'élément primitif}\label{sec:element-primitif}
\index{théorème!de l'élément primitif}
% ------------------------------------------------------------

\begin{theoreme}[Élément primitif]\label{thm:element-primitif}
Toute extension finie séparable $L/K$ est simple, c'est-à-dire de la forme
$L = K(\gamma)$ pour un certain $\gamma \in L$.
\end{theoreme}

\begin{proof}
\emph{Cas $K$ fini.} Si $K$ est fini, alors $L$ est fini. Le groupe
multiplicatif $L^*$ est cyclique (c'est un groupe fini d'un corps).
Soit $\gamma$ un générateur de $L^*$. Alors $L = K(\gamma)$.

\medskip\noindent
\emph{Cas $K$ infini.} Il suffit de traiter le cas $L = K(\alpha, \beta)$
(le cas général s'en déduit par récurrence). On cherche $\gamma$ de la forme
$\gamma = \alpha + c\beta$ avec $c \in K$.

Soient $f = \irr(\alpha, K)$ de racines $\alpha = \alpha_1, \ldots, \alpha_m$
et $g = \irr(\beta, K)$ de racines $\beta = \beta_1, \ldots, \beta_n$ dans
une clôture algébrique, toutes distinctes (séparabilité).

Pour $i \in \{1, \ldots, m\}$ et $j \in \{2, \ldots, n\}$ (on exclut
$j = 1$), l'équation $\alpha_i + c\beta_j = \alpha + c\beta$, soit
$c = \frac{\alpha - \alpha_i}{\beta_j - \beta}$, a au plus une solution
en $c$. Comme $K$ est infini, on peut choisir $c \in K$ évitant toutes
ces valeurs.

Posons $\gamma = \alpha + c\beta$ et montrons $K(\gamma) = K(\alpha,\beta)$.
L'inclusion $K(\gamma) \subseteq K(\alpha,\beta)$ est évidente. Pour
l'autre, considérons $h(X) = f(\gamma - cX) \in K(\gamma)[X]$.
On a $h(\beta) = f(\gamma - c\beta) = f(\alpha) = 0$, donc $\beta$ est une
racine commune à $g(X)$ et $h(X)$ dans $K(\gamma)[X]$.

Montrons que $\beta$ est la \emph{seule} racine commune. Si $\beta_j$
(avec $j \neq 1$) était aussi racine de $h$, on aurait
$f(\gamma - c\beta_j) = 0$, donc $\gamma - c\beta_j = \alpha_i$ pour un
certain $i$, soit $\alpha + c\beta - c\beta_j = \alpha_i$, c'est-à-dire
$c = (\alpha - \alpha_i)/(\beta_j - \beta)$, contredisant le choix de $c$.

Donc $\pgcd(g, h)$ dans $K(\gamma)[X]$ a $\beta$ comme unique racine.
Puisque $g$ est séparable, $\pgcd(g, h) = X - \beta$, ce qui implique
$\beta \in K(\gamma)$. Alors $\alpha = \gamma - c\beta \in K(\gamma)$
aussi, d'où $K(\alpha,\beta) \subseteq K(\gamma)$.
\end{proof}

\begin{corollaire}[Application aux extensions galoisiennes]\label{cor:element-primitif-galois}
Toute extension galoisienne finie $L/K$ est de la forme $L = K(\gamma)$
pour un $\gamma \in L$.
\end{corollaire}

\begin{proof}
Une extension galoisienne est séparable, donc le théorème~\ref{thm:element-primitif} s'applique.
\end{proof}

% ------------------------------------------------------------
\section{Exercices}\label{sec:exercices-chap12}
% ------------------------------------------------------------

\begin{exercice}[Correspondance pour $V_4$]\label{exo:correspondance-V4}
Soit $L = \Q(\sqrt{2},\sqrt{5})$.
\begin{enumerate}[label=(\alph*)]
  \item Montrer que $\Gal(L/\Q) \cong V_4$.
  \item Déterminer tous les sous-corps de $L$ contenant $\Q$.
  \item Dessiner les treillis (sous-groupes et sous-corps).
\end{enumerate}
\end{exercice}

\begin{exercice}[Correspondance pour $D_4$]\label{exo:correspondance-D4}
\index{groupe!diédral $D_4$}
Soit $L = \Q(\sqrt[4]{2}, i)$, corps de décomposition de $X^4 - 2$.
\begin{enumerate}[label=(\alph*)]
  \item Montrer que $[L:\Q] = 8$ et $\Gal(L/\Q) \cong D_4$.
  \item Déterminer tous les sous-groupes de $D_4$ (il y en a $10$).
  \item Pour chaque sous-groupe, identifier le sous-corps fixe correspondant.
  \item Quels sous-corps intermédiaires sont normaux sur $\Q$ ?
\end{enumerate}
\end{exercice}

\begin{exercice}[Groupe de Galois et discriminant]\label{exo:galois-discriminant}
\index{discriminant}
Soit $f \in K[X]$ un polynôme séparable de degré $n$ et $L$ son corps de
décomposition. Le groupe de Galois $G = \Gal(L/K)$ se plonge dans $S_n$.
Montrer que $G \subseteq A_n$ si et seulement si le discriminant
$\Delta(f)$ est un carré dans $K$.
\end{exercice}

\begin{exercice}[Théorème fondamental pour les corps finis]\label{exo:FTGT-corps-finis}
Donner la correspondance de Galois complète pour $\F_{p^{30}}/\F_p$.
Combien y a-t-il de sous-corps intermédiaires ?
\end{exercice}

\begin{exercice}[Sous-extensions d'une extension cyclotomique]\label{exo:sous-ext-cyclotomique}
Déterminer tous les sous-corps de $\Q(\zeta_7)$ contenant $\Q$.
\emph{Indication :} $\Gal(\Q(\zeta_7)/\Q) \cong (\Z/7\Z)^* \cong \Z/6\Z$.
\end{exercice}

\begin{exercice}[Groupe de Galois d'un polynôme de degré 4]\label{exo:galois-degre-4}
Déterminer le groupe de Galois de $f = X^4 - 2$ sur $\Q$, puis sur
$\Q(i)$, puis sur $\Q(\sqrt{2})$.
\end{exercice}

\begin{exercice}[Élément primitif explicite]\label{exo:element-primitif-explicite}
Trouver un élément primitif $\gamma$ tel que
$\Q(\sqrt{2},\sqrt{3}) = \Q(\gamma)$, et déterminer le polynôme minimal
de $\gamma$ sur $\Q$.
\end{exercice}

\begin{exercice}[Compositum et correspondance]\label{exo:compositum-correspondance}
Soient $E_1, E_2$ deux corps intermédiaires d'une extension galoisienne
$L/K$ de groupe $G$. Montrer que~:
\begin{enumerate}[label=(\alph*)]
  \item $\Gal(L/E_1 E_2) = \Gal(L/E_1) \cap \Gal(L/E_2)$,
  \item $\Gal(L/E_1 \cap E_2) = \langle \Gal(L/E_1), \Gal(L/E_2)\rangle$
    (le sous-groupe engendré).
\end{enumerate}
\end{exercice}

\begin{exercice}[Fermeture normale]\label{exo:fermeture-normale-correspondance}
Soit $L/K$ galoisienne de groupe $G$ et $E$ un corps intermédiaire avec
$H = \Gal(L/E)$. Montrer que la plus petite extension normale de $K$
contenue dans $L$ et contenant $E$ correspond au sous-groupe
$\bigcap_{\sigma \in G} \sigma H \sigma^{-1}$ (le plus grand sous-groupe
distingué de $G$ contenu dans $H$).
\end{exercice}

\begin{exercice}[Application du théorème de l'élément primitif]\label{exo:application-element-primitif}
Montrer qu'il n'existe qu'un nombre fini de sous-corps intermédiaires
pour toute extension galoisienne finie $L/K$.
\emph{Indication :} utiliser l'élément primitif et la correspondance de Galois.
\end{exercice}

\begin{exercice}[Correspondance de Galois en caractéristique $p$]\label{exo:galois-car-p}
Soit $K = \F_p(t)$ et $f = X^p - t \in K[X]$. Montrer que le corps
de décomposition de $f$ sur $K$ est une extension inséparable de degré $p$.
En déduire que $|\Aut(L/K)| < [L:K]$ et que la correspondance de Galois
ne s'applique pas.
\end{exercice}

\bigskip
\begin{center}
\fbox{\parbox{0.85\textwidth}{
\textbf{Résumé du chapitre~\ref{chap:theoreme-fondamental}.}
Le théorème fondamental de la théorie de Galois établit, pour une extension
galoisienne finie $L/K$ de groupe $G$, une bijection anti-isomorphe entre
sous-groupes de $G$ et corps intermédiaires, via $H \mapsto L^H$ et
$E \mapsto \Gal(L/E)$. Les degrés correspondent aux indices, et la normalité
d'une sous-extension correspond à la distinguabilité du sous-groupe associé.
Le théorème de l'élément primitif assure que toute extension séparable finie
est simple.
}}
\end{center}

% ============================================================
%  CHAPITRE 13 : APPLICATIONS
% ============================================================
\chapter{Applications de la théorie de Galois}\label{chap:applications}

\bigskip
\noindent
Ce chapitre illustre la puissance de la théorie de Galois à travers trois
domaines d'application classiques~: les constructions à la règle et au compas,
la résolution des équations polynomiales par radicaux, et le théorème
fondamental de l'algèbre. Nous verrons comment des questions géométriques et
algébriques vieilles de plus de deux millénaires trouvent leur réponse
définitive grâce à la machinerie développée dans les chapitres précédents.

% ============================================================
\section{Constructibilité à la règle et au compas}\label{sec:constructibilite}
\index{constructibilité!à la règle et au compas}
% ============================================================

\subsection{Nombres constructibles}\label{ssec:nombres-constructibles}

\begin{definition}[Construction élémentaire]\label{def:construction-elementaire}
\index{construction!élémentaire}
Partant d'un ensemble fini de points du plan $\R^2$ contenant au moins deux
points distincts, une \emph{étape de construction à la règle et au compas}
consiste en l'une des opérations suivantes~:
\begin{enumerate}[label=(\roman*)]
  \item Tracer la droite passant par deux points déjà construits.
  \item Tracer le cercle centré en un point construit et passant par un
    autre point construit.
  \item Marquer un point d'intersection de deux droites, de deux cercles,
    ou d'une droite et d'un cercle déjà tracés.
\end{enumerate}
\end{definition}

\begin{definition}[Nombre constructible]\label{def:nombre-constructible}
\index{nombre!constructible}
Un nombre réel $\alpha$ est dit \emph{constructible} si, partant des points
$(0,0)$ et $(1,0)$, on peut construire le point $(\alpha, 0)$ en un nombre
fini d'étapes à la règle et au compas.
\end{definition}

\begin{proposition}[Les nombres constructibles forment un corps]\label{prop:constructibles-corps}
L'ensemble $\mathcal{C}$ des nombres constructibles est un sous-corps de $\R$
contenant $\Q$. De plus, $\mathcal{C}$ est stable par extraction de racines
carrées~: si $\alpha \in \mathcal{C}$ et $\alpha > 0$, alors
$\sqrt{\alpha} \in \mathcal{C}$.
\end{proposition}

\begin{proof}
\emph{$\Q \subset \mathcal{C}$.} On construit d'abord tous les entiers par
report de longueurs. Puis les rationnels par le théorème de Thalès.

\emph{Stabilité par $+, -, \times, /$}. L'addition et la soustraction se
réalisent par report de longueurs sur une droite. La multiplication et la
division s'obtiennent par le théorème de Thalès~: si $a, b > 0$ sont
constructibles, on construit $ab$ en traçant deux droites sécantes, portant
les longueurs $1$ et $a$ sur l'une, $b$ sur l'autre, et en utilisant une
parallèle. La division est analogue.

\emph{Racine carrée.} Si $\alpha > 0$ est constructible, on trace un
segment de longueur $1 + \alpha$, on construit le cercle de diamètre
ce segment, et la hauteur issue du point de séparation a pour longueur
$\sqrt{\alpha}$ (par le théorème de la moyenne géométrique
dans un demi-cercle).
\end{proof}

\begin{center}
\begin{tikzpicture}[scale=1.5]
  % Construction de racine carrée
  \coordinate (A) at (0,0);
  \coordinate (B) at (1,0);
  \coordinate (C) at (3,0);
  % Segment [AC] de longueur 1+a, avec a=2 ici
  \draw[thick] (A) -- (C);
  \draw[thick,blue!60!black] (1.5,0) circle (1.5);
  \coordinate (H) at (1,{sqrt(2)});
  \draw[thick,red!60!black] (B) -- (H);
  \draw[dashed,gray] (A) -- (H) -- (C);
  \fill (A) circle (1pt) node[below] {$0$};
  \fill (B) circle (1pt) node[below] {$1$};
  \fill (C) circle (1pt) node[below] {$1+\alpha$};
  \fill (H) circle (1pt) node[above right] {$\sqrt{\alpha}$};
  \node[below,gray] at (0.5,0) {$1$};
  \node[below,gray] at (2,0) {$\alpha$};
\end{tikzpicture}
\end{center}

\begin{theoreme}[Critère de constructibilité]\label{thm:critere-constructibilite}
\index{constructibilité!critère de degré}
Si $\alpha \in \R$ est constructible, alors il existe une tour de corps
\[
  \Q = K_0 \subset K_1 \subset \cdots \subset K_r
\]
avec $\alpha \in K_r$ et $[K_{i+1}:K_i] = 2$ pour tout $i$. En particulier,
\[
  [\Q(\alpha):\Q] \text{ divise } 2^r, \quad\text{donc }
  [\Q(\alpha):\Q] \text{ est une puissance de } 2.
\]
\end{theoreme}

\begin{proof}
Chaque étape de construction ajoute un point obtenu comme intersection de
droites et cercles dont les coefficients sont dans le corps $K_i$ engendré
par les coordonnées des points déjà construits.

\emph{Intersection de deux droites.} Le système linéaire a ses coefficients
dans $K_i$, donc la solution est dans $K_i$. Pas d'extension nécessaire.

\emph{Intersection d'une droite et d'un cercle.} On substitue l'équation
de la droite dans celle du cercle, obtenant une équation du second degré à
coefficients dans $K_i$. Les coordonnées du point d'intersection sont dans
une extension de degré au plus $2$ de $K_i$.

\emph{Intersection de deux cercles.} En soustrayant les deux équations
$x^2 + y^2 + a_j x + b_j y + c_j = 0$ ($j = 1, 2$), on obtient une
équation linéaire, ce qui ramène au cas précédent.

Donc à chaque étape, on étend $K_i$ en $K_{i+1}$ avec
$[K_{i+1}:K_i] \in \{1, 2\}$. En ne conservant que les extensions
strictes, on obtient la tour voulue. La formule des degrés donne
$[K_r:\Q] = 2^r$ et $[\Q(\alpha):\Q]$ divise $[K_r:\Q] = 2^r$.
\end{proof}

\begin{corollaire}[Condition nécessaire]\label{cor:condition-necessaire-constr}
Si $\alpha \in \R$ est constructible, alors $[\Q(\alpha):\Q]$ est une
puissance de $2$.
\end{corollaire}

\subsection{Impossibilités classiques}\label{ssec:impossibilites}

\begin{theoreme}[Duplication du cube]\label{thm:duplication-cube}
\index{duplication du cube}
Il est impossible de doubler le cube à la règle et au compas~: le nombre
$\sqrt[3]{2}$ n'est pas constructible.
\end{theoreme}

\begin{proof}
Le polynôme $X^3 - 2$ est irréductible sur $\Q$ (critère d'Eisenstein avec
$p = 2$, ou test des racines rationnelles). Donc
$[\Q(\sqrt[3]{2}):\Q] = 3$. Or $3$ n'est pas une puissance de $2$. Par le
corollaire~\ref{cor:condition-necessaire-constr}, $\sqrt[3]{2}$ n'est pas
constructible.
\end{proof}

\begin{theoreme}[Trisection de l'angle]\label{thm:trisection-angle}
\index{trisection de l'angle}
Il est impossible de trisecter un angle arbitraire à la règle et au compas.
Plus précisément, l'angle de $60°$ n'est pas trisectable.
\end{theoreme}

\begin{proof}
Trisecter $60°$ revient à construire $\cos(20°)$. Or la formule
$\cos(3\theta) = 4\cos^3(\theta) - 3\cos(\theta)$ avec $\theta = 20°$
donne $\cos(60°) = 1/2$, donc $x = \cos(20°)$ vérifie
$8x^3 - 6x - 1 = 0$. Le changement de variable $y = 2x$ donne
$y^3 - 3y - 1 = 0$. Ce polynôme est irréductible sur $\Q$ (il n'a pas de
racine rationnelle, les candidats $\pm 1$ ne convenant pas). Donc
$[\Q(\cos 20°):\Q] = 3$, ce qui n'est pas une puissance de $2$.
\end{proof}

\begin{theoreme}[Quadrature du cercle]\label{thm:quadrature-cercle}
\index{quadrature du cercle}
Il est impossible de construire à la règle et au compas un carré de même aire
qu'un cercle donné. En d'autres termes, $\sqrt{\pi}$ (et donc $\pi$) n'est
pas constructible.
\end{theoreme}

\begin{proof}[Justification]
Lindemann a démontré en 1882 que $\pi$ est \emph{transcendant}\index{nombre!transcendant}
sur $\Q$, c'est-à-dire qu'il n'est racine d'aucun polynôme non nul à
coefficients rationnels. A fortiori, $[\Q(\pi):\Q]$ est infini, donc $\pi$
n'est pas constructible.
\end{proof}

\subsection{Constructibilité des polygones réguliers}\label{ssec:polygones-reguliers}
\index{polygone régulier!constructibilité}

\begin{theoreme}[Gauss--Wantzel]\label{thm:gauss-wantzel}
\index{théorème!de Gauss--Wantzel}
\index{Gauss--Wantzel}
Le polygone régulier à $n$ côtés ($n \geq 3$) est constructible à la règle et
au compas si et seulement si
\[
  n = 2^k \, p_1 \, p_2 \cdots p_s
\]
où $k \geq 0$ et les $p_i$ sont des nombres premiers de Fermat deux à deux
distincts\index{Fermat!nombre premier de}, c'est-à-dire de la forme
$p = 2^{2^m} + 1$.
\end{theoreme}

\begin{remarque}[Nombres premiers de Fermat]\label{rem:fermat-primes}
Les nombres premiers de Fermat connus sont~:
$F_0 = 3$, $F_1 = 5$, $F_2 = 17$, $F_3 = 257$, $F_4 = 65537$.
On ne sait pas s'il en existe d'autres. Euler a montré que
$F_5 = 4\,294\,967\,297 = 641 \times 6\,700\,417$ n'est pas premier.

Le théorème de Gauss--Wantzel implique par exemple que le polygone régulier
à $7$ côtés n'est pas constructible (car $7$ n'est pas un nombre premier de
Fermat), tandis que l'heptadécagone ($17$ côtés) l'est. La construction
explicite de l'heptadécagone par Gauss, alors âgé de $19$ ans, est l'une des
motivations qui l'ont poussé à choisir les mathématiques plutôt que la
philologie.
\end{remarque}

\begin{proof}[Idée de la preuve du théorème de Gauss--Wantzel]
Le polygone régulier à $n$ côtés est constructible si et seulement si
$\zeta_n = e^{2\pi i/n}$ est constructible, c'est-à-dire si et seulement si
$[\Q(\zeta_n):\Q] = \varphi(n)$ est une puissance de $2$. L'analyse de la
fonction $\varphi$ montre que $\varphi(n)$ est une puissance de $2$ si et
seulement si $n$ est de la forme $2^k p_1 \cdots p_s$ avec les $p_i$ des
premiers de Fermat distincts.
\end{proof}

% ============================================================
\section{Résolution par radicaux}\label{sec:resolution-radicaux}
\index{résolution!par radicaux}
% ============================================================

\subsection{Extensions radicales}\label{ssec:extensions-radicales}

\begin{definition}[Extension radicale]\label{def:extension-radicale}
\index{extension!radicale}
Une extension $L/K$ est dite \emph{radicale} s'il existe une tour
\[
  K = K_0 \subset K_1 \subset \cdots \subset K_r = L
\]
telle que pour chaque $i$, $K_{i+1} = K_i(\alpha_i)$ où $\alpha_i^{n_i}
\in K_i$ pour un certain entier $n_i \geq 1$. On dit que chaque étape est
une \emph{adjonction de radical}.
\end{definition}

\begin{definition}[Polynôme résoluble par radicaux]\label{def:resoluble-radicaux}
\index{polynôme!résoluble par radicaux}
Un polynôme $f \in K[X]$ est dit \emph{résoluble par radicaux} si son corps
de décomposition $L$ est contenu dans une extension radicale de $K$, c'est-à-dire
s'il existe une extension radicale $M/K$ avec $L \subseteq M$.
\end{definition}

\subsection{Groupes résolubles}\label{ssec:groupes-resolubles}

\begin{definition}[Groupe résoluble]\label{def:groupe-resoluble}
\index{groupe!résoluble}
Un groupe $G$ est dit \emph{résoluble} s'il existe une suite de sous-groupes
\[
  \{e\} = G_0 \trianglelefteq G_1 \trianglelefteq \cdots \trianglelefteq
  G_s = G
\]
telle que chaque quotient $G_{i+1}/G_i$ est abélien. Une telle suite est
appelée \emph{suite de composition abélienne}\index{suite de composition!abélienne}.
\end{definition}

\begin{exemple}[Groupes résolubles]\label{ex:groupes-resolubles}
\begin{enumerate}[label=(\roman*)]
  \item Tout groupe abélien est résoluble (prendre $G_0 = \{e\}$,
    $G_1 = G$).
  \item $S_3$ est résoluble~: $\{e\} \trianglelefteq A_3 \trianglelefteq
    S_3$ avec $A_3 \cong \Z/3\Z$ (abélien) et $S_3/A_3 \cong \Z/2\Z$
    (abélien).
  \item $S_4$ est résoluble~:
    $\{e\} \trianglelefteq V_4 \trianglelefteq A_4 \trianglelefteq S_4$
    avec $V_4 \cong (\Z/2\Z)^2$, $A_4/V_4 \cong \Z/3\Z$,
    $S_4/A_4 \cong \Z/2\Z$.
  \item Pour $n \leq 4$, $S_n$ est résoluble.
\end{enumerate}
\end{exemple}

\begin{theoreme}[$S_n$ n'est pas résoluble pour $n \geq 5$]\label{thm:Sn-non-resoluble}
\index{$S_n$!non résoluble}
Pour $n \geq 5$, le groupe symétrique $S_n$ n'est pas résoluble.
\end{theoreme}

\begin{proof}
Il suffit de montrer que $A_n$ est simple (c'est-à-dire n'a pas de
sous-groupe distingué non trivial propre) pour $n \geq 5$.

\emph{Fait clef :} $A_n$ est engendré par les $3$-cycles pour tout $n \geq 3$.
En effet, tout élément de $A_n$ est un produit d'un nombre pair de
transpositions, et $(ij)(jk) = (ijk)$, $(ij)(kl) = (ijk)(jkl)$.

\emph{Tout sous-groupe distingué non trivial de $A_n$ contient un $3$-cycle}
(pour $n \geq 5$). Soit $N \neq \{e\}$ un sous-groupe distingué de $A_n$.
On distingue des cas selon le type de cycle de $\sigma \in N \setminus \{e\}$~:

\begin{itemize}
  \item Si $\sigma$ contient un cycle de longueur $\geq 4$, disons
    $\sigma = (1\,2\,3\,4\,\ldots)\tau$, alors le commutateur
    $[\sigma, (1\,2\,3)] = \sigma(1\,2\,3)\sigma^{-1}(1\,3\,2) \in N$
    est un produit de $3$-cycles qui n'est pas l'identité, et un
    calcul montre qu'on peut extraire un $3$-cycle de $N$.
  \item Si $\sigma$ est un produit de transpositions disjointes,
    $\sigma = (1\,2)(3\,4)\tau$ avec $\tau$ fixant $\{1,2,3,4\}$,
    alors pour $n \geq 5$, le commutateur
    $[\sigma, (1\,2\,3)] = \sigma(1\,2\,3)\sigma^{-1}(1\,3\,2) = (1\,3)(2\,4)$
    est dans $N$. Puis $[(1\,3)(2\,4),\,(1\,2\,5)] \in N$ est un $3$-cycle.
  \item Si $\sigma$ est un $3$-cycle, c'est immédiat.
\end{itemize}

Donc $N$ contient un $3$-cycle. Puisque $N$ est distingué dans $A_n$ et que
tous les $3$-cycles sont conjugués dans $A_n$ (pour $n \geq 5$), $N$
contient tous les $3$-cycles, donc $N = A_n$.

Ainsi $A_n$ est simple pour $n \geq 5$. Si $S_n$ était résoluble, $A_n$
(en tant que sous-groupe d'un groupe résoluble) serait aussi résoluble. Mais un
groupe simple non abélien n'est pas résoluble (sa seule suite de composition
abélienne possible serait $\{e\} \trianglelefteq A_n$, et $A_n/\{e\} = A_n$
n'est pas abélien pour $n \geq 5$). Contradiction.
\end{proof}

\subsection{Le théorème de Galois}\label{ssec:theoreme-galois}

\begin{theoreme}[Galois]\label{thm:galois-resolubilite}
\index{théorème!de Galois (résolubilit\'e)}
Soit $K$ un corps de caractéristique $0$ et $f \in K[X]$ un polynôme
séparable. Soit $L$ le corps de décomposition de $f$ sur $K$. Alors~:
\[
  f \text{ est résoluble par radicaux} \iff \Gal(L/K) \text{ est résoluble}.
\]
\end{theoreme}

\begin{proof}[Esquisse de preuve]
($\Leftarrow$) Supposons $G = \Gal(L/K)$ résoluble. Quitte à adjoindre des
racines de l'unité appropriées, on peut supposer que $K$ contient les racines
$n$-ièmes de l'unité pour $n = |G|$. Alors la suite de composition
$\{e\} = G_0 \trianglelefteq G_1 \trianglelefteq \cdots \trianglelefteq
G_s = G$ avec quotients abéliens correspond, via la correspondance de Galois,
à une tour de corps
$L = E_0 \supset E_1 \supset \cdots \supset E_s = K$.
Chaque extension $E_{i-1}/E_i$ est galoisienne de groupe abélien $G_i/G_{i-1}$.
Par la théorie de Kummer\index{Kummer!théorie de}, une extension abélienne
d'un corps contenant les racines de l'unité appropriées est radicale.
Donc $L/K$ est contenu dans une extension radicale.

($\Rightarrow$) Si $f$ est résoluble par radicaux, il existe une extension
radicale $M/K$ contenant $L$. On peut supposer $M/K$ galoisienne (quitte à
prendre la clôture galoisienne). Chaque étape de la tour radicale correspond
à une extension cyclique (après adjonction de racines de l'unité), dont le
groupe de Galois est cyclique (donc abélien). Par un argument de raffinement
de la tour, on montre que $\Gal(M/K)$ est résoluble. Comme $\Gal(L/K)$ est
un quotient de $\Gal(M/K)$ (par le théorème fondamental), et qu'un quotient
d'un groupe résoluble est résoluble, on conclut.
\end{proof}

\subsection{Le théorème d'Abel--Ruffini}\label{ssec:abel-ruffini}

\begin{theoreme}[Abel--Ruffini]\label{thm:abel-ruffini}
\index{théorème!d'Abel--Ruffini}
\index{Abel--Ruffini}
L'\emph{équation générale} de degré $n \geq 5$ n'est pas résoluble par
radicaux. Plus précisément, il existe des polynômes de degré $5$ à
coefficients dans $\Q$ qui ne sont pas résolubles par radicaux.
\end{theoreme}

\begin{proof}
Par le théorème~\ref{thm:galois-resolubilite}, il suffit d'exhiber un
polynôme de degré $5$ dont le groupe de Galois sur $\Q$ est $S_5$. Or $S_5$
n'est pas résoluble (théorème~\ref{thm:Sn-non-resoluble}).
\end{proof}

\begin{exemple}[$X^5 - 4X + 2$ a pour groupe de Galois $S_5$]\label{ex:X5-4X+2}
\index{polynôme!non résoluble par radicaux}
Considérons $f = X^5 - 4X + 2 \in \Q[X]$.

\emph{Irréductibilité.} Par le critère d'Eisenstein appliqué avec $p = 2$,
$f$ est irréductible sur $\Q$.

\emph{Le groupe de Galois est $S_5$.} Le groupe $G = \Gal(L/\Q)$, où $L$ est
le corps de décomposition de $f$, est un sous-groupe transitif de $S_5$ (car
$f$ est irréductible). De plus~:
\begin{itemize}
  \item $f$ a exactement trois racines réelles et deux racines complexes
    conjuguées. La conjugaison complexe induit une transposition dans $G$.
  \item $f$ est irréductible de degré $5$, donc $5 \mid |G|$. Par le théorème
    de Cauchy, $G$ contient un élément d'ordre $5$, c'est-à-dire un $5$-cycle.
\end{itemize}
Un sous-groupe transitif de $S_5$ contenant une transposition et un
$5$-cycle est $S_5$ tout entier (résultat classique de théorie des groupes).
Donc $\Gal(L/\Q) = S_5$, et $f$ n'est pas résoluble par radicaux.
\end{exemple}

\begin{remarque}[Perspective historique]\label{rem:perspective-historique}
\index{Galois, Évariste}
\index{Abel, Niels Henrik}
L'impossibilité de la résolution par radicaux de l'équation générale de degré
$\geq 5$ fut d'abord prouvée par Abel (1824), mais sans le cadre conceptuel
permettant de traiter les polynômes \emph{particuliers}. C'est Galois, dans
un mémoire rédigé en 1831 (la veille de sa mort, dit la légende, bien que la
réalité soit plus nuancée), qui a fourni le critère définitif~: un polynôme
est résoluble par radicaux si et seulement si son «~groupe~» (ce que nous
appelons aujourd'hui le groupe de Galois) est résoluble. Les manuscrits de
Galois, d'abord incompris, furent publiés par Liouville en 1846 et
développés par Jordan, Dedekind et Klein dans la seconde moitié du
\textsc{xix}\textsuperscript{e} siècle.
\end{remarque}

% ============================================================
\section{Théorème fondamental de l'algèbre}\label{sec:TFA}
\index{théorème!fondamental de l'algèbre}
% ============================================================

\begin{theoreme}[Théorème fondamental de l'algèbre]\label{thm:TFA}
Le corps $\C$ des nombres complexes est algébriquement clos~: tout polynôme
non constant $f \in \C[X]$ possède au moins une racine dans $\C$.
\end{theoreme}

\begin{proof}[Démonstration par la théorie de Galois]
Il suffit de montrer que $\C$ n'a pas d'extension algébrique propre, ou de
manière équivalente, que tout polynôme irréductible de $\C[X]$ est de degré $1$.
La preuve utilise trois faits~:
\begin{enumerate}[label=(\alph*)]
  \item Tout polynôme de $\R[X]$ de degré impair a une racine réelle (par
    le théorème des valeurs intermédiaires).
  \item $[\C:\R] = 2$, donc tout polynôme de degré $2$ sur $\R$ a ses
    racines dans $\C$.
  \item (Théorie de Sylow) Un $2$-groupe fini non trivial possède un
    sous-groupe d'indice $2$.
\end{enumerate}

Soit $f \in \R[X]$ irréductible. Par (a), $\deg(f)$ est pair.
Soit $L$ le corps de décomposition de $f$ sur $\R$~; c'est une extension
galoisienne. Posons $G = \Gal(L/\R)$ et $|G| = 2^a m$ avec $m$ impair.

Soit $H$ un $2$-sous-groupe de Sylow de $G$ et $E = L^H$. Alors
$[E:\R] = [G:H] = m$ est impair. Par le théorème de l'élément primitif,
$E = \R(\gamma)$ pour un certain $\gamma$, et
$[\R(\gamma):\R] = \deg(\irr(\gamma,\R))$ est impair. Par (a), ce degré
doit être $1$ (tout polynôme irréductible sur $\R$ de degré impair est de
degré $1$), donc $m = 1$ et $|G| = 2^a$.

Ainsi $G$ est un $2$-groupe. Si $a \geq 2$, par (c), $G$ possède un
sous-groupe $H'$ d'indice $2$, et $E' = L^{H'}$ vérifie $[E':\R] = 2$,
donc $E' = \R(\alpha)$ avec $\alpha$ racine d'un polynôme de degré $2$ sur
$\R$. Par (b), $\alpha \in \C$, donc $E' \subseteq \C$. Comme
$[E':\R] = 2 = [\C:\R]$, on a $E' = \C$. Donc $\C \subseteq L$ et
$[L:\C] = 2^{a-1}$.

Si $a - 1 \geq 1$, on recommence~: $\Gal(L/\C)$ est un $2$-groupe d'ordre
$2^{a-1} \geq 2$, qui possède un sous-groupe d'indice $2$. Le corps fixe
correspondant $F$ vérifie $[F:\C] = 2$. Mais tout polynôme de degré $2$
sur $\C$ a ses racines dans $\C$ (par la formule $(-b \pm \sqrt{b^2 - 4ac})/2a$
et l'existence des racines carrées dans $\C$). Donc il n'existe pas
d'extension de degré $2$ de $\C$, contradiction.

Donc $a - 1 = 0$, c'est-à-dire $a = 1$ et $L = \C$. Ainsi tout polynôme
irréductible de $\R[X]$ a son corps de décomposition contenu dans $\C$, d'où
$\deg(f) \leq 2$. Et si $\deg(f) = 2$, ses racines sont dans $\C$ par (b).

Finalement, tout polynôme non constant de $\C[X]$ est scindé dans $\C$.
\end{proof}

% ============================================================
\section{Ouverture : prolongements et perspectives}\label{sec:perspectives}
% ============================================================

La théorie de Galois, telle que nous l'avons développée dans ces chapitres,
constitue le point de départ de vastes domaines des mathématiques
contemporaines.

\begin{itemize}
  \item \textbf{Géométrie algébrique}\index{géométrie algébrique}.
    Le groupe fondamental étale d'un schéma est l'analogue du groupe de Galois~:
    il classifie les revêtements finis étales, de la même manière que le groupe
    de Galois classifie les extensions séparables. La correspondance de Galois
    pour les corps se généralise en une équivalence de catégories entre
    revêtements étales et ensembles munis d'une action continue du groupe
    fondamental (Grothendieck, \emph{SGA 1}, 1961).

  \item \textbf{Théorie algébrique des nombres}\index{théorie algébrique des nombres}.
    Le groupe de Galois absolu $\Gal(\overline{\Q}/\Q)$ est un objet central
    de la théorie des nombres. Son étude, notamment via les représentations
    galoisiennes et le programme de Langlands, est l'un des domaines les plus
    actifs de la recherche mathématique actuelle.

  \item \textbf{Théorie de Galois inverse}\index{théorie de Galois!inverse}.
    Tout groupe fini est-il le groupe de Galois d'une extension de $\Q$~?
    Ce problème, toujours ouvert en général, est résolu positivement pour de
    nombreuses classes de groupes (groupes résolubles par Shafarevich, groupes
    symétriques et alternés, groupes sporadiques\ldots).

  \item \textbf{Théorie de Galois différentielle}\index{théorie de Galois!différentielle}.
    Par analogie avec les extensions de corps, on peut associer un «~groupe de
    Galois différentiel~» à une équation différentielle linéaire. Ce groupe
    (un groupe algébrique) encode la nature de ses solutions, et une
    correspondance de Galois analogue au théorème fondamental est valable
    (Picard--Vessiot).

  \item \textbf{Topologie}\index{topologie!et théorie de Galois}.
    La correspondance entre sous-groupes du groupe fondamental $\pi_1(X, x_0)$
    et revêtements connexes de $X$ est formellement analogue à la
    correspondance de Galois. Ce n'est pas une coïncidence~: les deux théories
    sont des instances du même formalisme catégorique.
\end{itemize}

% ============================================================
\section{Exercices}\label{sec:exercices-chap13}
% ============================================================

\begin{exercice}[Constructibilité de $\cos(2\pi/9)$]\label{exo:cos-2pi-9}
Montrer que $\cos(2\pi/9)$ n'est pas constructible à la règle et au compas.
\emph{Indication :} montrer que $[\Q(\cos(2\pi/9)):\Q] = 3$.
\end{exercice}

\begin{exercice}[Constructibilité du pentagone]\label{exo:pentagone}
\index{pentagone régulier}
Montrer que le pentagone régulier est constructible à la règle et au compas.
\emph{Indication :} $\Phi_5(X) = X^4 + X^3 + X^2 + X + 1$ et
$\varphi(5) = 4 = 2^2$.
\end{exercice}

\begin{exercice}[Polygone à $17$ côtés]\label{exo:heptadecagone}
\index{heptadécagone}
Montrer que l'heptadécagone (polygone régulier à $17$ côtés) est constructible.
\emph{Indication :} $17$ est un nombre premier de Fermat, $\varphi(17) = 16 = 2^4$.
\end{exercice}

\begin{exercice}[Résolubilit\'e de $S_4$]\label{exo:resolubilite-S4}
Détailler la suite de composition
$\{e\} \trianglelefteq V_4 \trianglelefteq A_4 \trianglelefteq S_4$
et vérifier que chaque quotient est abélien.
\end{exercice}

\begin{exercice}[Groupe de Galois de $X^4 - 2$]\label{exo:galois-X4-2}
Montrer que le groupe de Galois de $X^4 - 2$ sur $\Q$ est $D_4$ (groupe
diédral d'ordre $8$). En déduire que $X^4 - 2$ est résoluble par radicaux.
\end{exercice}

\begin{exercice}[Groupe de Galois de $X^5 - 2$]\label{exo:galois-X5-2}
Déterminer le groupe de Galois de $X^5 - 2$ sur $\Q$. Ce polynôme est-il
résoluble par radicaux ?
\emph{Indication :} le groupe est $\Z/5\Z \rtimes (\Z/5\Z)^*$, un groupe
métacyclique d'ordre $20$, qui est résoluble.
\end{exercice}

\begin{exercice}[Non-résolubilit\'e explicite]\label{exo:non-resolubilite}
Montrer que $f = X^5 - 6X + 3$ n'est pas résoluble par radicaux sur $\Q$.
\emph{Indication :} vérifier l'irréductibilité (Eisenstein), compter les
racines réelles, exhiber une transposition dans le groupe de Galois via la
conjugaison complexe, et un $5$-cycle via le théorème de Cauchy.
\end{exercice}

\begin{exercice}[Théorème fondamental de l'algèbre --- approche analytique]\label{exo:TFA-analytique}
Donner une preuve du théorème fondamental de l'algèbre utilisant le principe
du module minimum (si $f \in \C[X]$ est non constant et $|f|$ atteint son
minimum en $z_0$, alors $f(z_0) = 0$).
\end{exercice}

\bigskip
\begin{center}
\fbox{\parbox{0.85\textwidth}{
\textbf{Résumé du chapitre~\ref{chap:applications}.}
La théorie de Galois résout trois classes de problèmes classiques~:
(1)~les constructions à la règle et au compas sont gouvernées par les
extensions de degré $2^r$, ce qui interdit la duplication du cube, la
trisection de l'angle et la quadrature du cercle~;
(2)~un polynôme est résoluble par radicaux si et seulement si son groupe
de Galois est résoluble, ce qui interdit la résolution par radicaux de
l'équation générale de degré $\geq 5$ (Abel--Ruffini)~;
(3)~$\C$ est algébriquement clos (théorème fondamental de l'algèbre).
}}
\end{center}


% ============================================================
%  BIBLIOGRAPHIE
% ============================================================
\begin{thebibliography}{99}

\bibitem{DummitFoote}
\textsc{Dummit}, D.\,S. et \textsc{Foote}, R.\,M.,
\emph{Abstract Algebra}, 3\textsuperscript{e} éd.,
John Wiley \& Sons, 2004.
\index{Dummit--Foote}

\bibitem{Lang}
\textsc{Lang}, S.,
\emph{Algebra}, 3\textsuperscript{e} éd., Graduate Texts in Mathematics 211,
Springer, 2002.
\index{Lang, Serge}

\bibitem{Artin}
\textsc{Artin}, E.,
\emph{Galois Theory}, 2\textsuperscript{e} éd.,
Dover Publications, 1998.
Réédition de l'exposé original de 1942 à Notre Dame.
\index{Artin, Emil}

\bibitem{Jacobson}
\textsc{Jacobson}, N.,
\emph{Basic Algebra I \& II},
W.\,H.\,Freeman, 1985--1989.
\index{Jacobson, Nathan}

\bibitem{Bourbaki}
\textsc{Bourbaki}, N.,
\emph{Algèbre}, Chapitres 4 à 7,
Springer, 2007.
\index{Bourbaki, Nicolas}

\bibitem{Perrin}
\textsc{Perrin}, D.,
\emph{Algèbre --- Cours de mathématiques de première année},
Ellipses, 1996.
\index{Perrin, Daniel}

\bibitem{Stewart}
\textsc{Stewart}, I.,
\emph{Galois Theory}, 4\textsuperscript{e} éd.,
CRC Press, 2015.
\index{Stewart, Ian}

\end{thebibliography}

% ============================================================
\printindex
\end{document}
